\makeatletter
\newif\ifBookIndexing
\BookIndexingfalse
\newcommand{\BookBoldPage}[1]{\textbf{\hyperpage{#1}}}
\DeclareRobustCommand{\BookRecordIndex}[1]{%
  \ifBookIndexing
    \expandafter\index\expandafter{\csname BookIndex@#1\endcsname|hyperpage}%
  \fi}
\DeclareRobustCommand{\BookIndexDefinition}[1]{%
  \ifBookIndexing
    \expandafter\index\expandafter{\csname BookIndex@#1\endcsname|BookBoldPage}%
  \fi}
\DeclareRobustCommand{\BookTerm}[2]{%
  \BookRecordIndex{#1}%
  \begingroup\let\BookRecordIndex\@gobble #2\endgroup}
\AddToHook{cmd/@outputpage/before}{\let\BookRecordIndex\@gobble}
\pdfstringdefDisableCommands{%
  \def\BookTerm#1#2{#2}%
  \let\BookRecordIndex\@gobble
  \let\BookIndexDefinition\@gobble}
\expandafter\def\csname BookIndex@T0001\endcsname{fanchou@∞-范畴（\ensuremath{\infty}-category）}
\expandafter\let\csname BookOriginal@trInfinityCategory\endcsname\trInfinityCategory
\renewcommand{\trInfinityCategory}{\BookTerm{T0001}{\csname BookOriginal@trInfinityCategory\endcsname}}
\expandafter\def\csname BookIndex@T0002\endcsname{nifanchou@拟范畴（quasi-category）}
\expandafter\let\csname BookOriginal@trQuasiCategory\endcsname\trQuasiCategory
\renewcommand{\trQuasiCategory}{\BookTerm{T0002}{\csname BookOriginal@trQuasiCategory\endcsname}}
\expandafter\def\csname BookIndex@T0003\endcsname{yuzhou@∞-宇宙（\ensuremath{\infty}-cosmos）}
\expandafter\let\csname BookOriginal@trInfinityCosmos\endcsname\trInfinityCosmos
\renewcommand{\trInfinityCosmos}{\BookTerm{T0003}{\csname BookOriginal@trInfinityCosmos\endcsname}}
\expandafter\def\csname BookIndex@T0004\endcsname{fanchou@范畴（category）}
\expandafter\let\csname BookOriginal@trCategory\endcsname\trCategory
\renewcommand{\trCategory}{\BookTerm{T0004}{\csname BookOriginal@trCategory\endcsname}}
\expandafter\def\csname BookIndex@T0005\endcsname{hanzi@函子（functor）}
\expandafter\let\csname BookOriginal@trFunctor\endcsname\trFunctor
\renewcommand{\trFunctor}{\BookTerm{T0005}{\csname BookOriginal@trFunctor\endcsname}}
\expandafter\def\csname BookIndex@T0006\endcsname{ziranbianhuan@自然变换（natural transformation）}
\expandafter\let\csname BookOriginal@trNaturalTransformation\endcsname\trNaturalTransformation
\renewcommand{\trNaturalTransformation}{\BookTerm{T0006}{\csname BookOriginal@trNaturalTransformation\endcsname}}
\expandafter\def\csname BookIndex@T0007\endcsname{dengjia@等价（equivalence）}
\expandafter\let\csname BookOriginal@trEquivalence\endcsname\trEquivalence
\renewcommand{\trEquivalence}{\BookTerm{T0007}{\csname BookOriginal@trEquivalence\endcsname}}
\expandafter\def\csname BookIndex@T0008\endcsname{jixian@极限（limit）}
\expandafter\let\csname BookOriginal@trLimit\endcsname\trLimit
\renewcommand{\trLimit}{\BookTerm{T0008}{\csname BookOriginal@trLimit\endcsname}}
\expandafter\def\csname BookIndex@T0009\endcsname{yujixian@余极限（colimit）}
\expandafter\let\csname BookOriginal@trColimit\endcsname\trColimit
\renewcommand{\trColimit}{\BookTerm{T0009}{\csname BookOriginal@trColimit\endcsname}}
\expandafter\def\csname BookIndex@T0010\endcsname{zhui@锥（cone）}
\expandafter\let\csname BookOriginal@trCone\endcsname\trCone
\renewcommand{\trCone}{\BookTerm{T0010}{\csname BookOriginal@trCone\endcsname}}
\expandafter\def\csname BookIndex@T0011\endcsname{bansui@伴随（adjunction）}
\expandafter\let\csname BookOriginal@trAdjunction\endcsname\trAdjunction
\renewcommand{\trAdjunction}{\BookTerm{T0011}{\csname BookOriginal@trAdjunction\endcsname}}
\expandafter\def\csname BookIndex@T0012\endcsname{kankuozhang@Kan 扩张（Kan extension）}
\expandafter\let\csname BookOriginal@trKanExtension\endcsname\trKanExtension
\renewcommand{\trKanExtension}{\BookTerm{T0012}{\csname BookOriginal@trKanExtension\endcsname}}
\expandafter\def\csname BookIndex@T0013\endcsname{tisheng@提升（lifting）}
\expandafter\let\csname BookOriginal@trLifting\endcsname\trLifting
\renewcommand{\trLifting}{\BookTerm{T0013}{\csname BookOriginal@trLifting\endcsname}}
\expandafter\def\csname BookIndex@T0014\endcsname{kuozhang@扩张（extension）}
\expandafter\let\csname BookOriginal@trExtension\endcsname\trExtension
\renewcommand{\trExtension}{\BookTerm{T0014}{\csname BookOriginal@trExtension\endcsname}}
\expandafter\def\csname BookIndex@T0015\endcsname{xianweihua@纤维化（fibration）}
\expandafter\let\csname BookOriginal@trFibration\endcsname\trFibration
\renewcommand{\trFibration}{\BookTerm{T0015}{\csname BookOriginal@trFibration\endcsname}}
\expandafter\def\csname BookIndex@T0016\endcsname{dengjiaxianweihua@等价纤维化（isofibration）}
\expandafter\let\csname BookOriginal@trIsofibration\endcsname\trIsofibration
\renewcommand{\trIsofibration}{\BookTerm{T0016}{\csname BookOriginal@trIsofibration\endcsname}}
\expandafter\def\csname BookIndex@T0017\endcsname{dikaerxianweihua@笛卡尔纤维化（cartesian fibration）}
\expandafter\let\csname BookOriginal@trCartesianFibration\endcsname\trCartesianFibration
\renewcommand{\trCartesianFibration}{\BookTerm{T0017}{\csname BookOriginal@trCartesianFibration\endcsname}}
\expandafter\def\csname BookIndex@T0018\endcsname{shuangcexianweihua@双侧纤维化（two-sided fibration）}
\expandafter\let\csname BookOriginal@trTwoSidedFibration\endcsname\trTwoSidedFibration
\renewcommand{\trTwoSidedFibration}{\BookTerm{T0018}{\csname BookOriginal@trTwoSidedFibration\endcsname}}
\expandafter\def\csname BookIndex@T0019\endcsname{douhaofanchou@逗号范畴（comma category）}
\expandafter\let\csname BookOriginal@trCommaCategory\endcsname\trCommaCategory
\renewcommand{\trCommaCategory}{\BookTerm{T0019}{\csname BookOriginal@trCommaCategory\endcsname}}
\expandafter\def\csname BookIndex@T0020\endcsname{douhaofanchou@逗号∞-范畴（comma \ensuremath{\infty}-category）}
\expandafter\let\csname BookOriginal@trCommaInfinityCategory\endcsname\trCommaInfinityCategory
\renewcommand{\trCommaInfinityCategory}{\BookTerm{T0020}{\csname BookOriginal@trCommaInfinityCategory\endcsname}}
\expandafter\def\csname BookIndex@T0021\endcsname{fugaihanzi@覆盖函子（smothering functor）}
\expandafter\let\csname BookOriginal@trSmotheringFunctor\endcsname\trSmotheringFunctor
\renewcommand{\trSmotheringFunctor}{\BookTerm{T0021}{\csname BookOriginal@trSmotheringFunctor\endcsname}}
\expandafter\def\csname BookIndex@T0022\endcsname{mo@模（module）}
\expandafter\let\csname BookOriginal@trModule\endcsname\trModule
\renewcommand{\trModule}{\BookTerm{T0022}{\csname BookOriginal@trModule\endcsname}}
\expandafter\def\csname BookIndex@T0023\endcsname{xunizhuangbei@虚拟装备（virtual equipment）}
\expandafter\let\csname BookOriginal@trVirtualEquipment\endcsname\trVirtualEquipment
\renewcommand{\trVirtualEquipment}{\BookTerm{T0023}{\csname BookOriginal@trVirtualEquipment\endcsname}}
\expandafter\def\csname BookIndex@T0024\endcsname{moxingwuguanxing@模型无关性（model independence）}
\expandafter\let\csname BookOriginal@trModelIndependence\endcsname\trModelIndependence
\renewcommand{\trModelIndependence}{\BookTerm{T0024}{\csname BookOriginal@trModelIndependence\endcsname}}
\expandafter\def\csname BookIndex@T0025\endcsname{moxingfanchou@模型范畴（model category）}
\expandafter\let\csname BookOriginal@trModelCategory\endcsname\trModelCategory
\renewcommand{\trModelCategory}{\BookTerm{T0025}{\csname BookOriginal@trModelCategory\endcsname}}
\expandafter\def\csname BookIndex@T0026\endcsname{fujifanchou@富集范畴（enriched category）}
\expandafter\let\csname BookOriginal@trEnrichedCategory\endcsname\trEnrichedCategory
\renewcommand{\trEnrichedCategory}{\BookTerm{T0026}{\csname BookOriginal@trEnrichedCategory\endcsname}}
\expandafter\def\csname BookIndex@T0027\endcsname{danchunji@单纯集（simplicial set）}
\expandafter\let\csname BookOriginal@trSimplicialSet\endcsname\trSimplicialSet
\renewcommand{\trSimplicialSet}{\BookTerm{T0027}{\csname BookOriginal@trSimplicialSet\endcsname}}
\expandafter\def\csname BookIndex@T0028\endcsname{daibiaojidanchunji@带标记单纯集（marked simplicial set）}
\expandafter\let\csname BookOriginal@trMarkedSimplicialSet\endcsname\trMarkedSimplicialSet
\renewcommand{\trMarkedSimplicialSet}{\BookTerm{T0028}{\csname BookOriginal@trMarkedSimplicialSet\endcsname}}
\expandafter\def\csname BookIndex@T0029\endcsname{jiaquanjixian@加权极限（weighted limit）}
\expandafter\let\csname BookOriginal@trWeightedLimit\endcsname\trWeightedLimit
\renewcommand{\trWeightedLimit}{\BookTerm{T0029}{\csname BookOriginal@trWeightedLimit\endcsname}}
\expandafter\def\csname BookIndex@T0030\endcsname{yuzhoulunhanzi@宇宙论函子（cosmological functor）}
\expandafter\let\csname BookOriginal@trCosmologicalFunctor\endcsname\trCosmologicalFunctor
\renewcommand{\trCosmologicalFunctor}{\BookTerm{T0030}{\csname BookOriginal@trCosmologicalFunctor\endcsname}}
\expandafter\def\csname BookIndex@T0031\endcsname{shuangdengjia@双等价（biequivalence）}
\expandafter\let\csname BookOriginal@trBiequivalence\endcsname\trBiequivalence
\renewcommand{\trBiequivalence}{\BookTerm{T0031}{\csname BookOriginal@trBiequivalence\endcsname}}
\expandafter\def\csname BookIndex@T0032\endcsname{huanmoxing@换模型（change-of-model）}
\expandafter\let\csname BookOriginal@trChangeOfModel\endcsname\trChangeOfModel
\renewcommand{\trChangeOfModel}{\BookTerm{T0032}{\csname BookOriginal@trChangeOfModel\endcsname}}
\expandafter\def\csname BookIndex@T0033\endcsname{yonedayinli@Yoneda 引理（Yoneda lemma）}
\expandafter\let\csname BookOriginal@trYonedaLemma\endcsname\trYonedaLemma
\renewcommand{\trYonedaLemma}{\BookTerm{T0033}{\csname BookOriginal@trYonedaLemma\endcsname}}
\expandafter\def\csname BookIndex@T0034\endcsname{tonglun2fanchou@同伦 2-范畴（homotopy 2-category）}
\expandafter\let\csname BookOriginal@trHomotopyTwoCategory\endcsname\trHomotopyTwoCategory
\renewcommand{\trHomotopyTwoCategory}{\BookTerm{T0034}{\csname BookOriginal@trHomotopyTwoCategory\endcsname}}
\expandafter\def\csname BookIndex@T0035\endcsname{zhantietu@粘贴图（pasting diagram）}
\expandafter\let\csname BookOriginal@trPastingDiagram\endcsname\trPastingDiagram
\renewcommand{\trPastingDiagram}{\BookTerm{T0035}{\csname BookOriginal@trPastingDiagram\endcsname}}
\expandafter\def\csname BookIndex@T0036\endcsname{zhenghefangkuai@正合方块（exact square）}
\expandafter\let\csname BookOriginal@trExactSquare\endcsname\trExactSquare
\renewcommand{\trExactSquare}{\BookTerm{T0036}{\csname BookOriginal@trExactSquare\endcsname}}
\expandafter\def\csname BookIndex@T0037\endcsname{dikaerde@笛卡尔的（cartesian）}
\expandafter\let\csname BookOriginal@trCartesian\endcsname\trCartesian
\renewcommand{\trCartesian}{\BookTerm{T0037}{\csname BookOriginal@trCartesian\endcsname}}
\expandafter\def\csname BookIndex@T0038\endcsname{yudikaerde@余笛卡尔的（cocartesian）}
\expandafter\let\csname BookOriginal@trCocartesian\endcsname\trCocartesian
\renewcommand{\trCocartesian}{\BookTerm{T0038}{\csname BookOriginal@trCocartesian\endcsname}}
\expandafter\def\csname BookIndex@T0039\endcsname{lisande@离散的（discrete）}
\expandafter\let\csname BookOriginal@trDiscrete\endcsname\trDiscrete
\renewcommand{\trDiscrete}{\BookTerm{T0039}{\csname BookOriginal@trDiscrete\endcsname}}
\expandafter\def\csname BookIndex@T0040\endcsname{rouxingde@柔性的（flexible）}
\expandafter\let\csname BookOriginal@trFlexible\endcsname\trFlexible
\renewcommand{\trFlexible}{\BookTerm{T0040}{\csname BookOriginal@trFlexible\endcsname}}
\expandafter\def\csname BookIndex@T0041\endcsname{teyide@特异的（exotic）}
\expandafter\let\csname BookOriginal@trExotic\endcsname\trExotic
\renewcommand{\trExotic}{\BookTerm{T0041}{\csname BookOriginal@trExotic\endcsname}}
\expandafter\def\csname BookIndex@T0042\endcsname{xingshide@形式的（formal）}
\expandafter\let\csname BookOriginal@trFormal\endcsname\trFormal
\renewcommand{\trFormal}{\BookTerm{T0042}{\csname BookOriginal@trFormal\endcsname}}
\expandafter\def\csname BookIndex@T0043\endcsname{daidiande@带点的（pointed）}
\expandafter\let\csname BookOriginal@trPointed\endcsname\trPointed
\renewcommand{\trPointed}{\BookTerm{T0043}{\csname BookOriginal@trPointed\endcsname}}
\expandafter\def\csname BookIndex@T0044\endcsname{wendingde@稳定的（stable）}
\expandafter\let\csname BookOriginal@trStable\endcsname\trStable
\renewcommand{\trStable}{\BookTerm{T0044}{\csname BookOriginal@trStable\endcsname}}
\expandafter\def\csname BookIndex@T0045\endcsname{kebiaoshide@可表示的（representable）}
\expandafter\let\csname BookOriginal@trRepresentable\endcsname\trRepresentable
\renewcommand{\trRepresentable}{\BookTerm{T0045}{\csname BookOriginal@trRepresentable\endcsname}}
\expandafter\def\csname BookIndex@T0046\endcsname{shi@始（initial）}
\expandafter\let\csname BookOriginal@trInitial\endcsname\trInitial
\renewcommand{\trInitial}{\BookTerm{T0046}{\csname BookOriginal@trInitial\endcsname}}
\expandafter\def\csname BookIndex@T0047\endcsname{zhong@终（terminal）}
\expandafter\let\csname BookOriginal@trTerminal\endcsname\trTerminal
\renewcommand{\trTerminal}{\BookTerm{T0047}{\csname BookOriginal@trTerminal\endcsname}}
\expandafter\def\csname BookIndex@T0048\endcsname{xianweihuade@纤维化的（fibered）}
\expandafter\let\csname BookOriginal@trFibered\endcsname\trFibered
\renewcommand{\trFibered}{\BookTerm{T0048}{\csname BookOriginal@trFibered\endcsname}}
\expandafter\def\csname BookIndex@T0049\endcsname{fuzhuangde@复状的（complicial）}
\expandafter\let\csname BookOriginal@trComplicial\endcsname\trComplicial
\renewcommand{\trComplicial}{\BookTerm{T0049}{\csname BookOriginal@trComplicial\endcsname}}
\expandafter\def\csname BookIndex@T0050\endcsname{moxingwuguande@模型无关的（model independent）}
\expandafter\let\csname BookOriginal@trModelIndependent\endcsname\trModelIndependent
\renewcommand{\trModelIndependent}{\BookTerm{T0050}{\csname BookOriginal@trModelIndependent\endcsname}}
\expandafter\def\csname BookIndex@T0051\endcsname{xianweixingduixiangfanchou@纤维性对象范畴（category of fibrant objects）}
\expandafter\let\csname BookOriginal@trCategoryOfFibrantObjects\endcsname\trCategoryOfFibrantObjects
\renewcommand{\trCategoryOfFibrantObjects}{\BookTerm{T0051}{\csname BookOriginal@trCategoryOfFibrantObjects\endcsname}}
\expandafter\def\csname BookIndex@T0052\endcsname{ruofenjiexitong@弱分解系统（weak factorization system）}
\expandafter\let\csname BookOriginal@trWeakFactorizationSystem\endcsname\trWeakFactorizationSystem
\renewcommand{\trWeakFactorizationSystem}{\BookTerm{T0052}{\csname BookOriginal@trWeakFactorizationSystem\endcsname}}
\expandafter\def\csname BookIndex@T0053\endcsname{baoqiangfuxing@胞腔复形（cell complex）}
\expandafter\let\csname BookOriginal@trCellComplex\endcsname\trCellComplex
\renewcommand{\trCellComplex}{\BookTerm{T0053}{\csname BookOriginal@trCellComplex\endcsname}}
\expandafter\def\csname BookIndex@T0054\endcsname{nifanchou@逆范畴（inverse category）}
\expandafter\let\csname BookOriginal@trInverseCategory\endcsname\trInverseCategory
\renewcommand{\trInverseCategory}{\BookTerm{T0054}{\csname BookOriginal@trInverseCategory\endcsname}}
\expandafter\def\csname BookIndex@T0055\endcsname{xiuzheng@修正（modification）}
\expandafter\let\csname BookOriginal@trModification\endcsname\trModification
\renewcommand{\trModification}{\BookTerm{T0055}{\csname BookOriginal@trModification\endcsname}}
\expandafter\def\csname BookIndex@T0056\endcsname{duiying@对应（mate）}
\expandafter\let\csname BookOriginal@trMate\endcsname\trMate
\renewcommand{\trMate}{\BookTerm{T0056}{\csname BookOriginal@trMate\endcsname}}
\expandafter\def\csname BookIndex@T0057\endcsname{gonge@共轭（conjugate）}
\expandafter\let\csname BookOriginal@trConjugate\endcsname\trConjugate
\renewcommand{\trConjugate}{\BookTerm{T0057}{\csname BookOriginal@trConjugate\endcsname}}
\expandafter\def\csname BookIndex@T0058\endcsname{dikaerji@笛卡尔积（cartesian product）}
\expandafter\let\csname BookOriginal@trCartesianProduct\endcsname\trCartesianProduct
\renewcommand{\trCartesianProduct}{\BookTerm{T0058}{\csname BookOriginal@trCartesianProduct\endcsname}}
\expandafter\def\csname BookIndex@T0059\endcsname{zihanzi@自函子（endofunctor）}
\expandafter\let\csname BookOriginal@trEndofunctor\endcsname\trEndofunctor
\renewcommand{\trEndofunctor}{\BookTerm{T0059}{\csname BookOriginal@trEndofunctor\endcsname}}
\expandafter\def\csname BookIndex@T0060\endcsname{denghuazi@等化子（equalizer）}
\expandafter\let\csname BookOriginal@trEqualizer\endcsname\trEqualizer
\renewcommand{\trEqualizer}{\BookTerm{T0060}{\csname BookOriginal@trEqualizer\endcsname}}
\expandafter\def\csname BookIndex@T0061\endcsname{zhuixingjixian@锥形极限（conical limit）}
\expandafter\let\csname BookOriginal@trConicalLimit\endcsname\trConicalLimit
\renewcommand{\trConicalLimit}{\BookTerm{T0061}{\csname BookOriginal@trConicalLimit\endcsname}}
\expandafter\def\csname BookIndex@T0062\endcsname{zhuxianweide@逐纤维的（fiberwise）}
\expandafter\let\csname BookOriginal@trFiberwise\endcsname\trFiberwise
\renewcommand{\trFiberwise}{\BookTerm{T0062}{\csname BookOriginal@trFiberwise\endcsname}}
\expandafter\def\csname BookIndex@T0063\endcsname{xuzhuangfuhe@须状复合（whiskered composite）}
\expandafter\let\csname BookOriginal@trWhiskering\endcsname\trWhiskering
\renewcommand{\trWhiskering}{\BookTerm{T0063}{\csname BookOriginal@trWhiskering\endcsname}}
\expandafter\def\csname BookIndex@T0064\endcsname{jubuxiaode@局部小的（locally small）}
\expandafter\let\csname BookOriginal@trLocallySmall\endcsname\trLocallySmall
\renewcommand{\trLocallySmall}{\BookTerm{T0064}{\csname BookOriginal@trLocallySmall\endcsname}}
\expandafter\def\csname BookIndex@T0065\endcsname{tonglunyujixian@同伦余极限（homotopy colimit）}
\expandafter\let\csname BookOriginal@trHomotopyColimit\endcsname\trHomotopyColimit
\renewcommand{\trHomotopyColimit}{\BookTerm{T0065}{\csname BookOriginal@trHomotopyColimit\endcsname}}
\expandafter\def\csname BookIndex@T0066\endcsname{tonglunjixian@同伦极限（homotopy limit）}
\expandafter\let\csname BookOriginal@trHomotopyLimit\endcsname\trHomotopyLimit
\renewcommand{\trHomotopyLimit}{\BookTerm{T0066}{\csname BookOriginal@trHomotopyLimit\endcsname}}
\expandafter\def\csname BookIndex@T0067\endcsname{danchunxing@单纯形（simplex）}
\expandafter\let\csname BookOriginal@trSimplex\endcsname\trSimplex
\renewcommand{\trSimplex}{\BookTerm{T0067}{\csname BookOriginal@trSimplex\endcsname}}
\expandafter\def\csname BookIndex@T0068\endcsname{danchunxingfanchou@单纯形范畴（simplex category）}
\expandafter\let\csname BookOriginal@trSimplexCategory\endcsname\trSimplexCategory
\renewcommand{\trSimplexCategory}{\BookTerm{T0068}{\csname BookOriginal@trSimplexCategory\endcsname}}
\expandafter\def\csname BookIndex@T0069\endcsname{danchun@单纯（simplicial）}
\expandafter\let\csname BookOriginal@trSimplicial\endcsname\trSimplicial
\renewcommand{\trSimplicial}{\BookTerm{T0069}{\csname BookOriginal@trSimplicial\endcsname}}
\expandafter\def\csname BookIndex@T0070\endcsname{danchunfanchou@单纯范畴（simplicial category）}
\expandafter\let\csname BookOriginal@trSimplicialCategory\endcsname\trSimplicialCategory
\renewcommand{\trSimplicialCategory}{\BookTerm{T0070}{\csname BookOriginal@trSimplicialCategory\endcsname}}
\expandafter\def\csname BookIndex@T0071\endcsname{danchunhanzi@单纯函子（simplicial functor）}
\expandafter\let\csname BookOriginal@trSimplicialFunctor\endcsname\trSimplicialFunctor
\renewcommand{\trSimplicialFunctor}{\BookTerm{T0071}{\csname BookOriginal@trSimplicialFunctor\endcsname}}
\expandafter\def\csname BookIndex@T0072\endcsname{danchunfujifanchou@单纯富集范畴（simplicially enriched category）}
\expandafter\let\csname BookOriginal@trSimpliciallyEnrichedCategory\endcsname\trSimpliciallyEnrichedCategory
\renewcommand{\trSimpliciallyEnrichedCategory}{\BookTerm{T0072}{\csname BookOriginal@trSimpliciallyEnrichedCategory\endcsname}}
\expandafter\def\csname BookIndex@T0073\endcsname{danchunfujijixian@单纯富集极限（simplicially enriched limit）}
\expandafter\let\csname BookOriginal@trSimpliciallyEnrichedLimit\endcsname\trSimpliciallyEnrichedLimit
\renewcommand{\trSimpliciallyEnrichedLimit}{\BookTerm{T0073}{\csname BookOriginal@trSimpliciallyEnrichedLimit\endcsname}}
\expandafter\def\csname BookIndex@T0074\endcsname{yudanchunduixiang@余单纯对象（cosimplicial object）}
\expandafter\let\csname BookOriginal@trCosimplicialObject\endcsname\trCosimplicialObject
\renewcommand{\trCosimplicialObject}{\BookTerm{T0074}{\csname BookOriginal@trCosimplicialObject\endcsname}}
\expandafter\def\csname BookIndex@T0075\endcsname{miansuanzi@面算子（face operator）}
\expandafter\let\csname BookOriginal@trFaceOperator\endcsname\trFaceOperator
\renewcommand{\trFaceOperator}{\BookTerm{T0075}{\csname BookOriginal@trFaceOperator\endcsname}}
\expandafter\def\csname BookIndex@T0076\endcsname{tuihuasuanzi@退化算子（degeneracy operator）}
\expandafter\let\csname BookOriginal@trDegeneracyOperator\endcsname\trDegeneracyOperator
\renewcommand{\trDegeneracyOperator}{\BookTerm{T0076}{\csname BookOriginal@trDegeneracyOperator\endcsname}}
\expandafter\def\csname BookIndex@T0077\endcsname{shenjing@神经（nerve）}
\expandafter\let\csname BookOriginal@trNerve\endcsname\trNerve
\renewcommand{\trNerve}{\BookTerm{T0077}{\csname BookOriginal@trNerve\endcsname}}
\expandafter\def\csname BookIndex@T0078\endcsname{xushufanchou@序数范畴（ordinal category）}
\expandafter\let\csname BookOriginal@trOrdinalCategory\endcsname\trOrdinalCategory
\renewcommand{\trOrdinalCategory}{\BookTerm{T0078}{\csname BookOriginal@trOrdinalCategory\endcsname}}
\expandafter\def\csname BookIndex@T0079\endcsname{jizhu@脊柱（spine）}
\expandafter\let\csname BookOriginal@trSpine\endcsname\trSpine
\renewcommand{\trSpine}{\BookTerm{T0079}{\csname BookOriginal@trSpine\endcsname}}
\expandafter\def\csname BookIndex@T0080\endcsname{gujia@骨架（skeleton）}
\expandafter\let\csname BookOriginal@trSkeleton\endcsname\trSkeleton
\renewcommand{\trSkeleton}{\BookTerm{T0080}{\csname BookOriginal@trSkeleton\endcsname}}
\expandafter\def\csname BookIndex@T0081\endcsname{2yugujiade@2-余骨架的（2-coskeletal）}
\expandafter\let\csname BookOriginal@trTwoCoskeletal\endcsname\trTwoCoskeletal
\renewcommand{\trTwoCoskeletal}{\BookTerm{T0081}{\csname BookOriginal@trTwoCoskeletal\endcsname}}
\expandafter\def\csname BookIndex@T0082\endcsname{haojiao@号角（horn）}
\expandafter\let\csname BookOriginal@trHorn\endcsname\trHorn
\renewcommand{\trHorn}{\BookTerm{T0082}{\csname BookOriginal@trHorn\endcsname}}
\expandafter\def\csname BookIndex@T0083\endcsname{bianjieqiumian@边界球面（boundary sphere）}
\expandafter\let\csname BookOriginal@trBoundarySphere\endcsname\trBoundarySphere
\renewcommand{\trBoundarySphere}{\BookTerm{T0083}{\csname BookOriginal@trBoundarySphere\endcsname}}
\expandafter\def\csname BookIndex@T0084\endcsname{tongyu@同余（congruence）}
\expandafter\let\csname BookOriginal@trCongruence\endcsname\trCongruence
\renewcommand{\trCongruence}{\BookTerm{T0084}{\csname BookOriginal@trCongruence\endcsname}}
\expandafter\def\csname BookIndex@T0085\endcsname{zifanyouxiangtu@自反有向图（reflexive directed graph）}
\expandafter\let\csname BookOriginal@trReflexiveDirectedGraph\endcsname\trReflexiveDirectedGraph
\renewcommand{\trReflexiveDirectedGraph}{\BookTerm{T0085}{\csname BookOriginal@trReflexiveDirectedGraph\endcsname}}
\expandafter\def\csname BookIndex@T0086\endcsname{ziyoufanchou@自由范畴（free category）}
\expandafter\let\csname BookOriginal@trFreeCategory\endcsname\trFreeCategory
\renewcommand{\trFreeCategory}{\BookTerm{T0086}{\csname BookOriginal@trFreeCategory\endcsname}}
\expandafter\def\csname BookIndex@T0087\endcsname{tonglun@同伦（homotopy）}
\expandafter\let\csname BookOriginal@trHomotopy\endcsname\trHomotopy
\renewcommand{\trHomotopy}{\BookTerm{T0087}{\csname BookOriginal@trHomotopy\endcsname}}
\expandafter\def\csname BookIndex@T0088\endcsname{tonglunde@同伦的（homotopic）}
\expandafter\let\csname BookOriginal@trHomotopic\endcsname\trHomotopic
\renewcommand{\trHomotopic}{\BookTerm{T0088}{\csname BookOriginal@trHomotopic\endcsname}}
\expandafter\def\csname BookIndex@T0089\endcsname{tonglunguanxi@同伦关系（homotopy relation）}
\expandafter\let\csname BookOriginal@trHomotopyRelation\endcsname\trHomotopyRelation
\renewcommand{\trHomotopyRelation}{\BookTerm{T0089}{\csname BookOriginal@trHomotopyRelation\endcsname}}
\expandafter\def\csname BookIndex@T0090\endcsname{tonglunfanchou@同伦范畴（homotopy category）}
\expandafter\let\csname BookOriginal@trHomotopyCategory\endcsname\trHomotopyCategory
\renewcommand{\trHomotopyCategory}{\BookTerm{T0090}{\csname BookOriginal@trHomotopyCategory\endcsname}}
\expandafter\def\csname BookIndex@T0091\endcsname{tonglunxianggantonggou@同伦相干同构（homotopy coherent isomorphism）}
\expandafter\let\csname BookOriginal@trHomotopyCoherentIsomorphism\endcsname\trHomotopyCoherentIsomorphism
\renewcommand{\trHomotopyCoherentIsomorphism}{\BookTerm{T0091}{\csname BookOriginal@trHomotopyCoherentIsomorphism\endcsname}}
\expandafter\def\csname BookIndex@T0092\endcsname{tonglundengjia@同伦等价（homotopy equivalence）}
\expandafter\let\csname BookOriginal@trHomotopyEquivalence\endcsname\trHomotopyEquivalence
\renewcommand{\trHomotopyEquivalence}{\BookTerm{T0092}{\csname BookOriginal@trHomotopyEquivalence\endcsname}}
\expandafter\def\csname BookIndex@T0093\endcsname{kanfuxing@Kan 复形（Kan complex）}
\expandafter\let\csname BookOriginal@trKanComplex\endcsname\trKanComplex
\renewcommand{\trKanComplex}{\BookTerm{T0093}{\csname BookOriginal@trKanComplex\endcsname}}
\expandafter\def\csname BookIndex@T0094\endcsname{kanxianweihua@Kan 纤维化（Kan fibration）}
\expandafter\let\csname BookOriginal@trKanFibration\endcsname\trKanFibration
\renewcommand{\trKanFibration}{\BookTerm{T0094}{\csname BookOriginal@trKanFibration\endcsname}}
\expandafter\def\csname BookIndex@T0095\endcsname{qunpei@群胚（groupoid）}
\expandafter\let\csname BookOriginal@trGroupoid\endcsname\trGroupoid
\renewcommand{\trGroupoid}{\BookTerm{T0095}{\csname BookOriginal@trGroupoid\endcsname}}
\expandafter\def\csname BookIndex@T0096\endcsname{qunpei@∞-群胚（\ensuremath{\infty}-groupoid）}
\expandafter\let\csname BookOriginal@trInfinityGroupoid\endcsname\trInfinityGroupoid
\renewcommand{\trInfinityGroupoid}{\BookTerm{T0096}{\csname BookOriginal@trInfinityGroupoid\endcsname}}
\expandafter\def\csname BookIndex@T0097\endcsname{qunpeihe@∞-群胚核（\ensuremath{\infty}-groupoid core）}
\expandafter\let\csname BookOriginal@trInfinityGroupoidCore\endcsname\trInfinityGroupoidCore
\renewcommand{\trInfinityGroupoidCore}{\BookTerm{T0097}{\csname BookOriginal@trInfinityGroupoidCore\endcsname}}
\expandafter\def\csname BookIndex@T0098\endcsname{hehanzi@核函子（core functor）}
\expandafter\let\csname BookOriginal@trCoreFunctor\endcsname\trCoreFunctor
\renewcommand{\trCoreFunctor}{\BookTerm{T0098}{\csname BookOriginal@trCoreFunctor\endcsname}}
\expandafter\def\csname BookIndex@T0099\endcsname{tonggou@同构（isomorphism）}
\expandafter\let\csname BookOriginal@trIsomorphism\endcsname\trIsomorphism
\renewcommand{\trIsomorphism}{\BookTerm{T0099}{\csname BookOriginal@trIsomorphism\endcsname}}
\expandafter\def\csname BookIndex@T0100\endcsname{ziyoujiantou@自由箭头（free-living arrow）}
\expandafter\let\csname BookOriginal@trFreeLivingArrow\endcsname\trFreeLivingArrow
\renewcommand{\trFreeLivingArrow}{\BookTerm{T0100}{\csname BookOriginal@trFreeLivingArrow\endcsname}}
\expandafter\def\csname BookIndex@T0101\endcsname{ziyoutonggou@自由同构（free-living isomorphism）}
\expandafter\let\csname BookOriginal@trFreeLivingIsomorphism\endcsname\trFreeLivingIsomorphism
\renewcommand{\trFreeLivingIsomorphism}{\BookTerm{T0101}{\csname BookOriginal@trFreeLivingIsomorphism\endcsname}}
\expandafter\def\csname BookIndex@T0102\endcsname{jidazikanfuxing@极大子 Kan 复形（maximal sub Kan complex）}
\expandafter\let\csname BookOriginal@trMaximalSubKanComplex\endcsname\trMaximalSubKanComplex
\renewcommand{\trMaximalSubKanComplex}{\BookTerm{T0102}{\csname BookOriginal@trMaximalSubKanComplex\endcsname}}
\expandafter\def\csname BookIndex@T0103\endcsname{dantaishe@单态射（monomorphism）}
\expandafter\let\csname BookOriginal@trMonomorphism\endcsname\trMonomorphism
\renewcommand{\trMonomorphism}{\BookTerm{T0103}{\csname BookOriginal@trMonomorphism\endcsname}}
\expandafter\def\csname BookIndex@T0104\endcsname{mantaishe@满态射（epimorphism）}
\expandafter\let\csname BookOriginal@trEpimorphism\endcsname\trEpimorphism
\renewcommand{\trEpimorphism}{\BookTerm{T0104}{\csname BookOriginal@trEpimorphism\endcsname}}
\expandafter\def\csname BookIndex@T0105\endcsname{lahui@拉回（pullback）}
\expandafter\let\csname BookOriginal@trPullback\endcsname\trPullback
\renewcommand{\trPullback}{\BookTerm{T0105}{\csname BookOriginal@trPullback\endcsname}}
\expandafter\def\csname BookIndex@T0106\endcsname{tuichu@推出（pushout）}
\expandafter\let\csname BookOriginal@trPushout\endcsname\trPushout
\renewcommand{\trPushout}{\BookTerm{T0106}{\csname BookOriginal@trPushout\endcsname}}
\expandafter\def\csname BookIndex@T0107\endcsname{shousuo@收缩（retract）}
\expandafter\let\csname BookOriginal@trRetract\endcsname\trRetract
\renewcommand{\trRetract}{\BookTerm{T0107}{\csname BookOriginal@trRetract\endcsname}}
\expandafter\def\csname BookIndex@T0108\endcsname{lengzhu@棱柱（prism）}
\expandafter\let\csname BookOriginal@trPrism\endcsname\trPrism
\renewcommand{\trPrism}{\BookTerm{T0108}{\csname BookOriginal@trPrism\endcsname}}
\expandafter\def\csname BookIndex@T0109\endcsname{jiemian@截面（section）}
\expandafter\let\csname BookOriginal@trSection\endcsname\trSection
\renewcommand{\trSection}{\BookTerm{T0109}{\csname BookOriginal@trSection\endcsname}}
\expandafter\def\csname BookIndex@T0110\endcsname{qujian@区间（interval）}
\expandafter\let\csname BookOriginal@trInterval\endcsname\trInterval
\renewcommand{\trInterval}{\BookTerm{T0110}{\csname BookOriginal@trInterval\endcsname}}
\expandafter\def\csname BookIndex@T0111\endcsname{zhishulixiang@指数理想（exponential ideal）}
\expandafter\let\csname BookOriginal@trExponentialIdeal\endcsname\trExponentialIdeal
\renewcommand{\trExponentialIdeal}{\BookTerm{T0111}{\csname BookOriginal@trExponentialIdeal\endcsname}}
\expandafter\def\csname BookIndex@T0112\endcsname{lujingduixiang@路径对象（path object）}
\expandafter\let\csname BookOriginal@trPathObject\endcsname\trPathObject
\renewcommand{\trPathObject}{\BookTerm{T0112}{\csname BookOriginal@trPathObject\endcsname}}
\expandafter\def\csname BookIndex@T0113\endcsname{tianchongti@填充体（filler）}
\expandafter\let\csname BookOriginal@trFiller\endcsname\trFiller
\renewcommand{\trFiller}{\BookTerm{T0113}{\csname BookOriginal@trFiller\endcsname}}
\expandafter\def\csname BookIndex@T0114\endcsname{yiwanghanzi@遗忘函子（forgetful functor）}
\expandafter\let\csname BookOriginal@trForgetfulFunctor\endcsname\trForgetfulFunctor
\renewcommand{\trForgetfulFunctor}{\BookTerm{T0114}{\csname BookOriginal@trForgetfulFunctor\endcsname}}
\expandafter\def\csname BookIndex@T0115\endcsname{chuandibibao@传递闭包（transitive closure）}
\expandafter\let\csname BookOriginal@trTransitiveClosure\endcsname\trTransitiveClosure
\renewcommand{\trTransitiveClosure}{\BookTerm{T0115}{\csname BookOriginal@trTransitiveClosure\endcsname}}
\expandafter\def\csname BookIndex@T0116\endcsname{moxingjiegou@模型结构（model structure）}
\expandafter\let\csname BookOriginal@trModelStructure\endcsname\trModelStructure
\renewcommand{\trModelStructure}{\BookTerm{T0116}{\csname BookOriginal@trModelStructure\endcsname}}
\expandafter\def\csname BookIndex@T0117\endcsname{xianweixing@纤维性（fibrant）}
\expandafter\let\csname BookOriginal@trFibrant\endcsname\trFibrant
\renewcommand{\trFibrant}{\BookTerm{T0117}{\csname BookOriginal@trFibrant\endcsname}}
\expandafter\def\csname BookIndex@T0118\endcsname{yuxianweixing@余纤维性（cofibrant）}
\expandafter\let\csname BookOriginal@trCofibrant\endcsname\trCofibrant
\renewcommand{\trCofibrant}{\BookTerm{T0118}{\csname BookOriginal@trCofibrant\endcsname}}
\expandafter\def\csname BookIndex@T0119\endcsname{ruodengjia@弱等价（weak equivalence）}
\expandafter\let\csname BookOriginal@trWeakEquivalence\endcsname\trWeakEquivalence
\renewcommand{\trWeakEquivalence}{\BookTerm{T0119}{\csname BookOriginal@trWeakEquivalence\endcsname}}
\expandafter\def\csname BookIndex@T0120\endcsname{pingfanxianweihua@平凡纤维化（trivial fibration）}
\expandafter\let\csname BookOriginal@trTrivialFibration\endcsname\trTrivialFibration
\renewcommand{\trTrivialFibration}{\BookTerm{T0120}{\csname BookOriginal@trTrivialFibration\endcsname}}
\expandafter\def\csname BookIndex@T0121\endcsname{keliemantaishe@可裂满态射（split epimorphism）}
\expandafter\let\csname BookOriginal@trSplitEpimorphism\endcsname\trSplitEpimorphism
\renewcommand{\trSplitEpimorphism}{\BookTerm{T0121}{\csname BookOriginal@trSplitEpimorphism\endcsname}}
\expandafter\def\csname BookIndex@T0122\endcsname{keliexianweitonglundengjia@可裂纤维同伦等价（split fiber homotopy equivalence）}
\expandafter\let\csname BookOriginal@trSplitFiberHomotopyEquivalence\endcsname\trSplitFiberHomotopyEquivalence
\renewcommand{\trSplitFiberHomotopyEquivalence}{\BookTerm{T0122}{\csname BookOriginal@trSplitFiberHomotopyEquivalence\endcsname}}
\expandafter\def\csname BookIndex@T0123\endcsname{benzhimande@本质满的（essentially surjective）}
\expandafter\let\csname BookOriginal@trEssentiallySurjective\endcsname\trEssentiallySurjective
\renewcommand{\trEssentiallySurjective}{\BookTerm{T0123}{\csname BookOriginal@trEssentiallySurjective\endcsname}}
\expandafter\def\csname BookIndex@T0124\endcsname{jubuquanzhongshide@局部全忠实的（locally fully faithful）}
\expandafter\let\csname BookOriginal@trLocallyFullyFaithful\endcsname\trLocallyFullyFaithful
\renewcommand{\trLocallyFullyFaithful}{\BookTerm{T0124}{\csname BookOriginal@trLocallyFullyFaithful\endcsname}}
\expandafter\def\csname BookIndex@T0125\endcsname{quanzhongshide@全忠实的（fully faithful）}
\expandafter\let\csname BookOriginal@trFullyFaithful\endcsname\trFullyFaithful
\renewcommand{\trFullyFaithful}{\BookTerm{T0125}{\csname BookOriginal@trFullyFaithful\endcsname}}
\expandafter\def\csname BookIndex@T0126\endcsname{dikaerfengbide@笛卡尔封闭的（cartesian closed）}
\expandafter\let\csname BookOriginal@trCartesianClosed\endcsname\trCartesianClosed
\renewcommand{\trCartesianClosed}{\BookTerm{T0126}{\csname BookOriginal@trCartesianClosed\endcsname}}
\expandafter\def\csname BookIndex@T0127\endcsname{yuwanbei@余完备（cocomplete）}
\expandafter\let\csname BookOriginal@trCocomplete\endcsname\trCocomplete
\renewcommand{\trCocomplete}{\BookTerm{T0127}{\csname BookOriginal@trCocomplete\endcsname}}
\expandafter\def\csname BookIndex@T0128\endcsname{huanji@换基（change of base）}
\expandafter\let\csname BookOriginal@trChangeOfBase\endcsname\trChangeOfBase
\renewcommand{\trChangeOfBase}{\BookTerm{T0128}{\csname BookOriginal@trChangeOfBase\endcsname}}
\expandafter\def\csname BookIndex@T0129\endcsname{yuzhoulunshuangdengjia@宇宙论双等价（cosmological biequivalence）}
\expandafter\let\csname BookOriginal@trCosmologicalBiequivalence\endcsname\trCosmologicalBiequivalence
\renewcommand{\trCosmologicalBiequivalence}{\BookTerm{T0129}{\csname BookOriginal@trCosmologicalBiequivalence\endcsname}}
\expandafter\def\csname BookIndex@T0130\endcsname{yuzhoulunqianru@宇宙论嵌入（cosmological embedding）}
\expandafter\let\csname BookOriginal@trCosmologicalEmbedding\endcsname\trCosmologicalEmbedding
\renewcommand{\trCosmologicalEmbedding}{\BookTerm{T0130}{\csname BookOriginal@trCosmologicalEmbedding\endcsname}}
\expandafter\def\csname BookIndex@T0131\endcsname{yuzhoulunjixiangainian@宇宙论极限概念（cosmological limit notions）}
\expandafter\let\csname BookOriginal@trCosmologicalLimitNotions\endcsname\trCosmologicalLimitNotions
\renewcommand{\trCosmologicalLimitNotions}{\BookTerm{T0131}{\csname BookOriginal@trCosmologicalLimitNotions\endcsname}}
\expandafter\def\csname BookIndex@T0132\endcsname{jubudengjia@局部等价（local equivalence）}
\expandafter\let\csname BookOriginal@trLocalEquivalence\endcsname\trLocalEquivalence
\renewcommand{\trLocalEquivalence}{\BookTerm{T0132}{\csname BookOriginal@trLocalEquivalence\endcsname}}
\expandafter\def\csname BookIndex@T0133\endcsname{manshedengjia@满射等价（surjective equivalence）}
\expandafter\let\csname BookOriginal@trSurjectiveEquivalence\endcsname\trSurjectiveEquivalence
\renewcommand{\trSurjectiveEquivalence}{\BookTerm{T0133}{\csname BookOriginal@trSurjectiveEquivalence\endcsname}}
\expandafter\def\csname BookIndex@T0134\endcsname{dinifanchou@底拟范畴（underlying quasi-category）}
\expandafter\let\csname BookOriginal@trUnderlyingQuasiCategory\endcsname\trUnderlyingQuasiCategory
\renewcommand{\trUnderlyingQuasiCategory}{\BookTerm{T0134}{\csname BookOriginal@trUnderlyingQuasiCategory\endcsname}}
\expandafter\def\csname BookIndex@T0135\endcsname{hanzikongjian@函子空间（functor space）}
\expandafter\let\csname BookOriginal@trFunctorSpace\endcsname\trFunctorSpace
\renewcommand{\trFunctorSpace}{\BookTerm{T0135}{\csname BookOriginal@trFunctorSpace\endcsname}}
\expandafter\def\csname BookIndex@T0136\endcsname{yuzhangliang@余张量（cotensor）}
\expandafter\let\csname BookOriginal@trCotensor\endcsname\trCotensor
\renewcommand{\trCotensor}{\BookTerm{T0136}{\csname BookOriginal@trCotensor\endcsname}}
\expandafter\def\csname BookIndex@T0137\endcsname{neibuhom@内部 hom（internal hom）}
\expandafter\let\csname BookOriginal@trInternalHom\endcsname\trInternalHom
\renewcommand{\trInternalHom}{\BookTerm{T0137}{\csname BookOriginal@trInternalHom\endcsname}}
\expandafter\def\csname BookIndex@T0138\endcsname{duiouyuzhou@对偶∞-宇宙（dual \ensuremath{\infty}-cosmos）}
\expandafter\let\csname BookOriginal@trDualInfinityCosmos\endcsname\trDualInfinityCosmos
\renewcommand{\trDualInfinityCosmos}{\BookTerm{T0138}{\csname BookOriginal@trDualInfinityCosmos\endcsname}}
\expandafter\def\csname BookIndex@T0139\endcsname{yuduiouyuzhou@余对偶∞-宇宙（co-dual \ensuremath{\infty}-cosmos）}
\expandafter\let\csname BookOriginal@trCodualInfinityCosmos\endcsname\trCodualInfinityCosmos
\renewcommand{\trCodualInfinityCosmos}{\BookTerm{T0139}{\csname BookOriginal@trCodualInfinityCosmos\endcsname}}
\expandafter\def\csname BookIndex@T0140\endcsname{fandanchunji@反单纯集（opposite simplicial set）}
\expandafter\let\csname BookOriginal@trOppositeSimplicialSet\endcsname\trOppositeSimplicialSet
\renewcommand{\trOppositeSimplicialSet}{\BookTerm{T0140}{\csname BookOriginal@trOppositeSimplicialSet\endcsname}}
\expandafter\def\csname BookIndex@T0141\endcsname{wanquansegalkongjian@完全 Segal 空间（complete Segal spaces）}
\expandafter\let\csname BookOriginal@trCompleteSegalSpaces\endcsname\trCompleteSegalSpaces
\renewcommand{\trCompleteSegalSpaces}{\BookTerm{T0141}{\csname BookOriginal@trCompleteSegalSpaces\endcsname}}
\expandafter\def\csname BookIndex@T0142\endcsname{segalfanchou@Segal 范畴（Segal categories）}
\expandafter\let\csname BookOriginal@trSegalCategories\endcsname\trSegalCategories
\renewcommand{\trSegalCategories}{\BookTerm{T0142}{\csname BookOriginal@trSegalCategories\endcsname}}
\expandafter\def\csname BookIndex@T0143\endcsname{1fuzhuangdanchunji@1-复状单纯集（1-complicial sets）}
\expandafter\let\csname BookOriginal@trOneComplicialSets\endcsname\trOneComplicialSets
\renewcommand{\trOneComplicialSets}{\BookTerm{T0143}{\csname BookOriginal@trOneComplicialSets\endcsname}}
\expandafter\def\csname BookIndex@T0144\endcsname{sanquerxingzhi@三取二性质（2-of-3 property）}
\expandafter\let\csname BookOriginal@trTwoOfThreeProperty\endcsname\trTwoOfThreeProperty
\renewcommand{\trTwoOfThreeProperty}{\BookTerm{T0144}{\csname BookOriginal@trTwoOfThreeProperty\endcsname}}
\expandafter\def\csname BookIndex@T0145\endcsname{sanquliuxingzhi@三取六性质（2-of-6 property）}
\expandafter\let\csname BookOriginal@trTwoOfSixProperty\endcsname\trTwoOfSixProperty
\renewcommand{\trTwoOfSixProperty}{\BookTerm{T0145}{\csname BookOriginal@trTwoOfSixProperty\endcsname}}
\expandafter\def\csname BookIndex@T0146\endcsname{yonedaqianru@Yoneda 嵌入（Yoneda embedding）}
\expandafter\let\csname BookOriginal@trYonedaEmbedding\endcsname\trYonedaEmbedding
\renewcommand{\trYonedaEmbedding}{\BookTerm{T0146}{\csname BookOriginal@trYonedaEmbedding\endcsname}}
\expandafter\def\csname BookIndex@T0147\endcsname{baoqiang@胞腔（cell）}
\expandafter\let\csname BookOriginal@trCell\endcsname\trCell
\renewcommand{\trCell}{\BookTerm{T0147}{\csname BookOriginal@trCell\endcsname}}
\expandafter\def\csname BookIndex@T0148\endcsname{2hanzi@2-函子（2-functor）}
\expandafter\let\csname BookOriginal@trTwoFunctor\endcsname\trTwoFunctor
\renewcommand{\trTwoFunctor}{\BookTerm{T0148}{\csname BookOriginal@trTwoFunctor\endcsname}}
\expandafter\def\csname BookIndex@T0149\endcsname{ziranbianhuan@∞-自然变换（\ensuremath{\infty}-natural transformation）}
\expandafter\let\csname BookOriginal@trInfinityNaturalTransformation\endcsname\trInfinityNaturalTransformation
\renewcommand{\trInfinityNaturalTransformation}{\BookTerm{T0149}{\csname BookOriginal@trInfinityNaturalTransformation\endcsname}}
\expandafter\def\csname BookIndex@T0150\endcsname{zirantonggou@∞-自然同构（\ensuremath{\infty}-natural isomorphism）}
\expandafter\let\csname BookOriginal@trInfinityNaturalIsomorphism\endcsname\trInfinityNaturalIsomorphism
\renewcommand{\trInfinityNaturalIsomorphism}{\BookTerm{T0150}{\csname BookOriginal@trInfinityNaturalIsomorphism\endcsname}}
\expandafter\def\csname BookIndex@T0151\endcsname{2zhongduixiang@2-终对象（2-terminal object）}
\expandafter\let\csname BookOriginal@trTwoTerminalObject\endcsname\trTwoTerminalObject
\renewcommand{\trTwoTerminalObject}{\BookTerm{T0151}{\csname BookOriginal@trTwoTerminalObject\endcsname}}
\expandafter\def\csname BookIndex@T0152\endcsname{yingshenifanchou@映射拟范畴（mapping quasi-category）}
\expandafter\let\csname BookOriginal@trMappingQuasiCategory\endcsname\trMappingQuasiCategory
\renewcommand{\trMappingQuasiCategory}{\BookTerm{T0152}{\csname BookOriginal@trMappingQuasiCategory\endcsname}}
\expandafter\def\csname BookIndex@T0153\endcsname{kenide@可逆的（invertible）}
\expandafter\let\csname BookOriginal@trInvertible\endcsname\trInvertible
\renewcommand{\trInvertible}{\BookTerm{T0153}{\csname BookOriginal@trInvertible\endcsname}}
\expandafter\def\csname BookIndex@T0154\endcsname{bianyuan@边缘（boundary）}
\expandafter\let\csname BookOriginal@trBoundary\endcsname\trBoundary
\renewcommand{\trBoundary}{\BookTerm{T0154}{\csname BookOriginal@trBoundary\endcsname}}
\expandafter\def\csname BookIndex@T0155\endcsname{houfuhe@后复合（post-composition）}
\expandafter\let\csname BookOriginal@trPostComposition\endcsname\trPostComposition
\renewcommand{\trPostComposition}{\BookTerm{T0155}{\csname BookOriginal@trPostComposition\endcsname}}
\expandafter\def\csname BookIndex@T0156\endcsname{zuobansui@左伴随（left adjoint）}
\expandafter\let\csname BookOriginal@trLeftAdjoint\endcsname\trLeftAdjoint
\renewcommand{\trLeftAdjoint}{\BookTerm{T0156}{\csname BookOriginal@trLeftAdjoint\endcsname}}
\expandafter\def\csname BookIndex@T0157\endcsname{youbansui@右伴随（right adjoint）}
\expandafter\let\csname BookOriginal@trRightAdjoint\endcsname\trRightAdjoint
\renewcommand{\trRightAdjoint}{\BookTerm{T0157}{\csname BookOriginal@trRightAdjoint\endcsname}}
\expandafter\def\csname BookIndex@T0158\endcsname{danwei@单位（unit）}
\expandafter\let\csname BookOriginal@trUnit\endcsname\trUnit
\renewcommand{\trUnit}{\BookTerm{T0158}{\csname BookOriginal@trUnit\endcsname}}
\expandafter\def\csname BookIndex@T0159\endcsname{yudanwei@余单位（counit）}
\expandafter\let\csname BookOriginal@trCounit\endcsname\trCounit
\renewcommand{\trCounit}{\BookTerm{T0159}{\csname BookOriginal@trCounit\endcsname}}
\expandafter\def\csname BookIndex@T0160\endcsname{sanjiaodengshi@三角等式（triangle equalities）}
\expandafter\let\csname BookOriginal@trTriangleEqualities\endcsname\trTriangleEqualities
\renewcommand{\trTriangleEqualities}{\BookTerm{T0160}{\csname BookOriginal@trTriangleEqualities\endcsname}}
\expandafter\def\csname BookIndex@T0161\endcsname{bansuidengjia@伴随等价（adjoint equivalence）}
\expandafter\let\csname BookOriginal@trAdjointEquivalence\endcsname\trAdjointEquivalence
\renewcommand{\trAdjointEquivalence}{\BookTerm{T0161}{\csname BookOriginal@trAdjointEquivalence\endcsname}}
\expandafter\def\csname BookIndex@T0162\endcsname{tonglunxiangganbansui@同伦相干伴随（homotopy coherent adjunction）}
\expandafter\let\csname BookOriginal@trHomotopyCoherentAdjunction\endcsname\trHomotopyCoherentAdjunction
\renewcommand{\trHomotopyCoherentAdjunction}{\BookTerm{T0162}{\csname BookOriginal@trHomotopyCoherentAdjunction\endcsname}}
\expandafter\def\csname BookIndex@T0163\endcsname{bansuisanyuanzu@伴随三元组（adjoint triple）}
\expandafter\let\csname BookOriginal@trAdjointTriple\endcsname\trAdjointTriple
\renewcommand{\trAdjointTriple}{\BookTerm{T0163}{\csname BookOriginal@trAdjointTriple\endcsname}}
\expandafter\def\csname BookIndex@T0164\endcsname{guangyiyuansu@广义元素（generalized element）}
\expandafter\let\csname BookOriginal@trGeneralizedElement\endcsname\trGeneralizedElement
\renewcommand{\trGeneralizedElement}{\BookTerm{T0164}{\csname BookOriginal@trGeneralizedElement\endcsname}}
\expandafter\def\csname BookIndex@T0165\endcsname{yuansu@元素（element）}
\expandafter\let\csname BookOriginal@trElement\endcsname\trElement
\renewcommand{\trElement}{\BookTerm{T0165}{\csname BookOriginal@trElement\endcsname}}
\expandafter\def\csname BookIndex@T0166\endcsname{fanxingzhi@泛性质（universal property）}
\expandafter\let\csname BookOriginal@trUniversalProperty\endcsname\trUniversalProperty
\renewcommand{\trUniversalProperty}{\BookTerm{T0166}{\csname BookOriginal@trUniversalProperty\endcsname}}
\expandafter\def\csname BookIndex@T0167\endcsname{kebiaoshideshide@可表示地始的（representably initial）}
\expandafter\let\csname BookOriginal@trRepresentablyInitial\endcsname\trRepresentablyInitial
\renewcommand{\trRepresentablyInitial}{\BookTerm{T0167}{\csname BookOriginal@trRepresentablyInitial\endcsname}}
\expandafter\def\csname BookIndex@T0168\endcsname{tonglunshide@同伦始的（homotopy initial）}
\expandafter\let\csname BookOriginal@trHomotopyInitial\endcsname\trHomotopyInitial
\renewcommand{\trHomotopyInitial}{\BookTerm{T0168}{\csname BookOriginal@trHomotopyInitial\endcsname}}
\expandafter\def\csname BookIndex@T0169\endcsname{wanbeide@完备的（complete）}
\expandafter\let\csname BookOriginal@trComplete\endcsname\trComplete
\renewcommand{\trComplete}{\BookTerm{T0169}{\csname BookOriginal@trComplete\endcsname}}
\expandafter\def\csname BookIndex@T0170\endcsname{changzhihanzi@常值函子（constant functor）}
\expandafter\let\csname BookOriginal@trConstantFunctor\endcsname\trConstantFunctor
\renewcommand{\trConstantFunctor}{\BookTerm{T0170}{\csname BookOriginal@trConstantFunctor\endcsname}}
\expandafter\def\csname BookIndex@T0171\endcsname{midengde@幂等的（idempotent）}
\expandafter\let\csname BookOriginal@trIdempotent\endcsname\trIdempotent
\renewcommand{\trIdempotent}{\BookTerm{T0171}{\csname BookOriginal@trIdempotent\endcsname}}
\expandafter\def\csname BookIndex@T0172\endcsname{youni@右逆（right inverse）}
\expandafter\let\csname BookOriginal@trRightInverse\endcsname\trRightInverse
\renewcommand{\trRightInverse}{\BookTerm{T0172}{\csname BookOriginal@trRightInverse\endcsname}}
\expandafter\def\csname BookIndex@T0173\endcsname{fanyuansu@泛元素（generic element）}
\expandafter\let\csname BookOriginal@trGenericElement\endcsname\trGenericElement
\renewcommand{\trGenericElement}{\BookTerm{T0173}{\csname BookOriginal@trGenericElement\endcsname}}
\expandafter\def\csname BookIndex@T0174\endcsname{tubiao@图表（diagram）}
\expandafter\let\csname BookOriginal@trDiagram\endcsname\trDiagram
\renewcommand{\trDiagram}{\BookTerm{T0174}{\csname BookOriginal@trDiagram\endcsname}}
\expandafter\def\csname BookIndex@T0175\endcsname{kongtubiao@空图表（empty diagram）}
\expandafter\let\csname BookOriginal@trEmptyDiagram\endcsname\trEmptyDiagram
\renewcommand{\trEmptyDiagram}{\BookTerm{T0175}{\csname BookOriginal@trEmptyDiagram\endcsname}}
\expandafter\def\csname BookIndex@T0176\endcsname{ji@积（product）}
\expandafter\let\csname BookOriginal@trProduct\endcsname\trProduct
\renewcommand{\trProduct}{\BookTerm{T0176}{\csname BookOriginal@trProduct\endcsname}}
\expandafter\def\csname BookIndex@T0177\endcsname{yuji@余积（coproduct）}
\expandafter\let\csname BookOriginal@trCoproduct\endcsname\trCoproduct
\renewcommand{\trCoproduct}{\BookTerm{T0177}{\csname BookOriginal@trCoproduct\endcsname}}
\expandafter\def\csname BookIndex@T0178\endcsname{danchunduixiang@单纯对象（simplicial object）}
\expandafter\let\csname BookOriginal@trSimplicialObject\endcsname\trSimplicialObject
\renewcommand{\trSimplicialObject}{\BookTerm{T0178}{\csname BookOriginal@trSimplicialObject\endcsname}}
\expandafter\def\csname BookIndex@T0179\endcsname{kelie@可裂（split）}
\expandafter\let\csname BookOriginal@trSplit\endcsname\trSplit
\renewcommand{\trSplit}{\BookTerm{T0179}{\csname BookOriginal@trSplit\endcsname}}
\expandafter\def\csname BookIndex@T0180\endcsname{zengguangde@增广的（augmented）}
\expandafter\let\csname BookOriginal@trAugmented\endcsname\trAugmented
\renewcommand{\trAugmented}{\BookTerm{T0180}{\csname BookOriginal@trAugmented\endcsname}}
\expandafter\def\csname BookIndex@T0181\endcsname{jiheshixian@几何实现（geometric realization）}
\expandafter\let\csname BookOriginal@trGeometricRealization\endcsname\trGeometricRealization
\renewcommand{\trGeometricRealization}{\BookTerm{T0181}{\csname BookOriginal@trGeometricRealization\endcsname}}
\expandafter\def\csname BookIndex@T0182\endcsname{quanhua@全化（totalization）}
\expandafter\let\csname BookOriginal@trTotalization\endcsname\trTotalization
\renewcommand{\trTotalization}{\BookTerm{T0182}{\csname BookOriginal@trTotalization\endcsname}}
\expandafter\def\csname BookIndex@T0183\endcsname{fugaide@覆盖的（smothering）}
\expandafter\let\csname BookOriginal@trSmothering\endcsname\trSmothering
\renewcommand{\trSmothering}{\BookTerm{T0183}{\csname BookOriginal@trSmothering\endcsname}}
\expandafter\def\csname BookIndex@T0184\endcsname{baoshoude@保守的（conservative）}
\expandafter\let\csname BookOriginal@trConservative\endcsname\trConservative
\renewcommand{\trConservative}{\BookTerm{T0184}{\csname BookOriginal@trConservative\endcsname}}
\expandafter\def\csname BookIndex@T0185\endcsname{xingzoujiantou@行走箭头（walking arrow）}
\expandafter\let\csname BookOriginal@trWalkingArrow\endcsname\trWalkingArrow
\renewcommand{\trWalkingArrow}{\BookTerm{T0185}{\csname BookOriginal@trWalkingArrow\endcsname}}
\expandafter\def\csname BookIndex@T0186\endcsname{jiantoufanchou@箭头范畴（category of arrows）}
\expandafter\let\csname BookOriginal@trCategoryOfArrows\endcsname\trCategoryOfArrows
\renewcommand{\trCategoryOfArrows}{\BookTerm{T0186}{\csname BookOriginal@trCategoryOfArrows\endcsname}}
\expandafter\def\csname BookIndex@T0187\endcsname{yukuadu@余跨度（cospan）}
\expandafter\let\csname BookOriginal@trCospan\endcsname\trCospan
\renewcommand{\trCospan}{\BookTerm{T0187}{\csname BookOriginal@trCospan\endcsname}}
\expandafter\def\csname BookIndex@T0188\endcsname{xianwei@纤维（fiber）}
\expandafter\let\csname BookOriginal@trFiber\endcsname\trFiber
\renewcommand{\trFiber}{\BookTerm{T0188}{\csname BookOriginal@trFiber\endcsname}}
\expandafter\def\csname BookIndex@T0189\endcsname{ziyoude@自由的（free）}
\expandafter\let\csname BookOriginal@trFree\endcsname\trFree
\renewcommand{\trFree}{\BookTerm{T0189}{\csname BookOriginal@trFree\endcsname}}
\expandafter\def\csname BookIndex@T0190\endcsname{zhongshide@忠实的（faithful）}
\expandafter\let\csname BookOriginal@trFaithful\endcsname\trFaithful
\renewcommand{\trFaithful}{\BookTerm{T0190}{\csname BookOriginal@trFaithful\endcsname}}
\expandafter\def\csname BookIndex@T0191\endcsname{mande@满的（full）}
\expandafter\let\csname BookOriginal@trFull\endcsname\trFull
\renewcommand{\trFull}{\BookTerm{T0191}{\csname BookOriginal@trFull\endcsname}}
\expandafter\def\csname BookIndex@T0192\endcsname{jiantoufanchou@箭头∞-范畴（arrow \ensuremath{\infty}-category）}
\expandafter\let\csname BookOriginal@trArrowInfinityCategory\endcsname\trArrowInfinityCategory
\renewcommand{\trArrowInfinityCategory}{\BookTerm{T0192}{\csname BookOriginal@trArrowInfinityCategory\endcsname}}
\expandafter\def\csname BookIndex@T0193\endcsname{fanjiantou@泛箭头（generic arrow）}
\expandafter\let\csname BookOriginal@trGenericArrow\endcsname\trGenericArrow
\renewcommand{\trGenericArrow}{\BookTerm{T0193}{\csname BookOriginal@trGenericArrow\endcsname}}
\expandafter\def\csname BookIndex@T0194\endcsname{duandianqiuzhi@端点求值（endpoint evaluation）}
\expandafter\let\csname BookOriginal@trEndpointEvaluation\endcsname\trEndpointEvaluation
\renewcommand{\trEndpointEvaluation}{\BookTerm{T0194}{\csname BookOriginal@trEndpointEvaluation\endcsname}}
\expandafter\def\csname BookIndex@T0195\endcsname{duandianbaohan@端点包含（endpoint inclusion）}
\expandafter\let\csname BookOriginal@trEndpointInclusion\endcsname\trEndpointInclusion
\renewcommand{\trEndpointInclusion}{\BookTerm{T0195}{\csname BookOriginal@trEndpointInclusion\endcsname}}
\expandafter\def\csname BookIndex@T0196\endcsname{dingyiyuqiuzhi@定义域求值（domain evaluation）}
\expandafter\let\csname BookOriginal@trDomainEvaluation\endcsname\trDomainEvaluation
\renewcommand{\trDomainEvaluation}{\BookTerm{T0196}{\csname BookOriginal@trDomainEvaluation\endcsname}}
\expandafter\def\csname BookIndex@T0197\endcsname{yudingyiyuqiuzhi@余定义域求值（codomain evaluation）}
\expandafter\let\csname BookOriginal@trCodomainEvaluation\endcsname\trCodomainEvaluation
\renewcommand{\trCodomainEvaluation}{\BookTerm{T0197}{\csname BookOriginal@trCodomainEvaluation\endcsname}}
\expandafter\def\csname BookIndex@T0198\endcsname{xianweihuadengjia@纤维化等价（fibered equivalence）}
\expandafter\let\csname BookOriginal@trFiberedEquivalence\endcsname\trFiberedEquivalence
\renewcommand{\trFiberedEquivalence}{\BookTerm{T0198}{\csname BookOriginal@trFiberedEquivalence\endcsname}}
\expandafter\def\csname BookIndex@T0199\endcsname{xianweihuatonggou@纤维化同构（fibered isomorphism）}
\expandafter\let\csname BookOriginal@trFiberedIsomorphism\endcsname\trFiberedIsomorphism
\renewcommand{\trFiberedIsomorphism}{\BookTerm{T0199}{\csname BookOriginal@trFiberedIsomorphism\endcsname}}
\expandafter\def\csname BookIndex@T0200\endcsname{xianweihuanidengjia@纤维化逆等价（fibered inverse equivalence）}
\expandafter\let\csname BookOriginal@trFiberedInverseEquivalence\endcsname\trFiberedInverseEquivalence
\renewcommand{\trFiberedInverseEquivalence}{\BookTerm{T0200}{\csname BookOriginal@trFiberedInverseEquivalence\endcsname}}
\expandafter\def\csname BookIndex@T0201\endcsname{qiepianyuzhou@切片∞-宇宙（slice \ensuremath{\infty}-cosmos）}
\expandafter\let\csname BookOriginal@trSliceInfinityCosmos\endcsname\trSliceInfinityCosmos
\renewcommand{\trSliceInfinityCosmos}{\BookTerm{T0201}{\csname BookOriginal@trSliceInfinityCosmos\endcsname}}
\expandafter\def\csname BookIndex@T0202\endcsname{1baoqiangguina@1-胞腔归纳（1-cell induction）}
\expandafter\let\csname BookOriginal@trOneCellInduction\endcsname\trOneCellInduction
\renewcommand{\trOneCellInduction}{\BookTerm{T0202}{\csname BookOriginal@trOneCellInduction\endcsname}}
\expandafter\def\csname BookIndex@T0203\endcsname{2baoqiangguina@2-胞腔归纳（2-cell induction）}
\expandafter\let\csname BookOriginal@trTwoCellInduction\endcsname\trTwoCellInduction
\renewcommand{\trTwoCellInduction}{\BookTerm{T0203}{\csname BookOriginal@trTwoCellInduction\endcsname}}
\expandafter\def\csname BookIndex@T0204\endcsname{2baoqiangbaoshouxing@2-胞腔保守性（2-cell conservativity）}
\expandafter\let\csname BookOriginal@trTwoCellConservativity\endcsname\trTwoCellConservativity
\renewcommand{\trTwoCellConservativity}{\BookTerm{T0204}{\csname BookOriginal@trTwoCellConservativity\endcsname}}
\expandafter\def\csname BookIndex@T0205\endcsname{homfanchou@hom-范畴（hom category）}
\expandafter\let\csname BookOriginal@trHomCategory\endcsname\trHomCategory
\renewcommand{\trHomCategory}{\BookTerm{T0205}{\csname BookOriginal@trHomCategory\endcsname}}
\expandafter\def\csname BookIndex@T0206\endcsname{ruodikaerde@弱笛卡尔的（weakly cartesian）}
\expandafter\let\csname BookOriginal@trWeaklyCartesian\endcsname\trWeaklyCartesian
\renewcommand{\trWeaklyCartesian}{\BookTerm{T0206}{\csname BookOriginal@trWeaklyCartesian\endcsname}}
\expandafter\def\csname BookIndex@T0207\endcsname{ruodikaerfangkuai@弱笛卡尔方块（weakly cartesian square）}
\expandafter\let\csname BookOriginal@trWeaklyCartesianSquare\endcsname\trWeaklyCartesianSquare
\renewcommand{\trWeaklyCartesianSquare}{\BookTerm{T0207}{\csname BookOriginal@trWeaklyCartesianSquare\endcsname}}
\expandafter\def\csname BookIndex@T0208\endcsname{youguyoude@右固有的（right proper）}
\expandafter\let\csname BookOriginal@trRightProper\endcsname\trRightProper
\renewcommand{\trRightProper}{\BookTerm{T0208}{\csname BookOriginal@trRightProper\endcsname}}
\expandafter\def\csname BookIndex@T0209\endcsname{youguyouxing@右固有性（right properness）}
\expandafter\let\csname BookOriginal@trRightProperness\endcsname\trRightProperness
\renewcommand{\trRightProperness}{\BookTerm{T0209}{\csname BookOriginal@trRightProperness\endcsname}}
\expandafter\def\csname BookIndex@T0210\endcsname{douhaozhui@逗号锥（comma cone）}
\expandafter\let\csname BookOriginal@trCommaCone\endcsname\trCommaCone
\renewcommand{\trCommaCone}{\BookTerm{T0210}{\csname BookOriginal@trCommaCone\endcsname}}
\expandafter\def\csname BookIndex@T0211\endcsname{yingshekongjian@映射空间（mapping space）}
\expandafter\let\csname BookOriginal@trMappingSpace\endcsname\trMappingSpace
\renewcommand{\trMappingSpace}{\BookTerm{T0211}{\csname BookOriginal@trMappingSpace\endcsname}}
\expandafter\def\csname BookIndex@T0212\endcsname{kuadu@跨度（span）}
\expandafter\let\csname BookOriginal@trSpan\endcsname\trSpan
\renewcommand{\trSpan}{\BookTerm{T0212}{\csname BookOriginal@trSpan\endcsname}}
\expandafter\def\csname BookIndex@T0213\endcsname{fansongchide@反松弛的（oplax）}
\expandafter\let\csname BookOriginal@trOplax\endcsname\trOplax
\renewcommand{\trOplax}{\BookTerm{T0213}{\csname BookOriginal@trOplax\endcsname}}
\expandafter\def\csname BookIndex@T0214\endcsname{leibnizzhangliang@Leibniz 张量（Leibniz tensor）}
\expandafter\let\csname BookOriginal@trLeibnizTensor\endcsname\trLeibnizTensor
\renewcommand{\trLeibnizTensor}{\BookTerm{T0214}{\csname BookOriginal@trLeibnizTensor\endcsname}}
\expandafter\def\csname BookIndex@T0215\endcsname{leibnizyuzhangliang@Leibniz 余张量（Leibniz cotensor）}
\expandafter\let\csname BookOriginal@trLeibnizCotensor\endcsname\trLeibnizCotensor
\renewcommand{\trLeibnizCotensor}{\BookTerm{T0215}{\csname BookOriginal@trLeibnizCotensor\endcsname}}
\expandafter\def\csname BookIndex@T0216\endcsname{jiaozheng@矫正（rectify）}
\expandafter\let\csname BookOriginal@trRectify\endcsname\trRectify
\renewcommand{\trRectify}{\BookTerm{T0216}{\csname BookOriginal@trRectify\endcsname}}
\expandafter\def\csname BookIndex@T0217\endcsname{nijiaohuande@拟交换的（pseudo commutative）}
\expandafter\let\csname BookOriginal@trPseudoCommutative\endcsname\trPseudoCommutative
\renewcommand{\trPseudoCommutative}{\BookTerm{T0217}{\csname BookOriginal@trPseudoCommutative\endcsname}}
\expandafter\def\csname BookIndex@T0218\endcsname{zhuifanchou@锥∞-范畴（\ensuremath{\infty}-category of cones）}
\expandafter\let\csname BookOriginal@trInfinityCategoryOfCones\endcsname\trInfinityCategoryOfCones
\renewcommand{\trInfinityCategoryOfCones}{\BookTerm{T0218}{\csname BookOriginal@trInfinityCategoryOfCones\endcsname}}
\expandafter\def\csname BookIndex@T0219\endcsname{zhangliang@张量（tensor）}
\expandafter\let\csname BookOriginal@trTensor\endcsname\trTensor
\renewcommand{\trTensor}{\BookTerm{T0219}{\csname BookOriginal@trTensor\endcsname}}
\expandafter\def\csname BookIndex@T0220\endcsname{gudi@谷底（nadir）}
\expandafter\let\csname BookOriginal@trNadir\endcsname\trNadir
\renewcommand{\trNadir}{\BookTerm{T0220}{\csname BookOriginal@trNadir\endcsname}}
\expandafter\def\csname BookIndex@T0221\endcsname{lingyuansu@零元素（zero element）}
\expandafter\let\csname BookOriginal@trZeroElement\endcsname\trZeroElement
\renewcommand{\trZeroElement}{\BookTerm{T0221}{\csname BookOriginal@trZeroElement\endcsname}}
\expandafter\def\csname BookIndex@T0222\endcsname{huilu@回路（loops）}
\expandafter\let\csname BookOriginal@trLoops\endcsname\trLoops
\renewcommand{\trLoops}{\BookTerm{T0222}{\csname BookOriginal@trLoops\endcsname}}
\expandafter\def\csname BookIndex@T0223\endcsname{xuangua@悬挂（suspension）}
\expandafter\let\csname BookOriginal@trSuspension\endcsname\trSuspension
\renewcommand{\trSuspension}{\BookTerm{T0223}{\csname BookOriginal@trSuspension\endcsname}}
\expandafter\def\csname BookIndex@T0224\endcsname{ziduixiang@子对象（subobject）}
\expandafter\let\csname BookOriginal@trSubobject\endcsname\trSubobject
\renewcommand{\trSubobject}{\BookTerm{T0224}{\csname BookOriginal@trSubobject\endcsname}}
\expandafter\def\csname BookIndex@T0225\endcsname{kesuode@可缩的（contractible）}
\expandafter\let\csname BookOriginal@trContractible\endcsname\trContractible
\renewcommand{\trContractible}{\BookTerm{T0225}{\csname BookOriginal@trContractible\endcsname}}
\expandafter\def\csname BookIndex@T0226\endcsname{linglun@零伦（nullhomotopy）}
\expandafter\let\csname BookOriginal@trNullhomotopy\endcsname\trNullhomotopy
\renewcommand{\trNullhomotopy}{\BookTerm{T0226}{\csname BookOriginal@trNullhomotopy\endcsname}}
\expandafter\def\csname BookIndex@T0227\endcsname{erhanzi@二函子（bifunctor）}
\expandafter\let\csname BookOriginal@trBifunctor\endcsname\trBifunctor
\renewcommand{\trBifunctor}{\BookTerm{T0227}{\csname BookOriginal@trBifunctor\endcsname}}
\expandafter\def\csname BookIndex@T0228\endcsname{shuangxianxingde@双线性的（bilinear）}
\expandafter\let\csname BookOriginal@trBilinear\endcsname\trBilinear
\renewcommand{\trBilinear}{\BookTerm{T0228}{\csname BookOriginal@trBilinear\endcsname}}
\expandafter\def\csname BookIndex@T0229\endcsname{abeierqun@阿贝尔群（abelian group）}
\expandafter\let\csname BookOriginal@trAbelianGroup\endcsname\trAbelianGroup
\renewcommand{\trAbelianGroup}{\BookTerm{T0229}{\csname BookOriginal@trAbelianGroup\endcsname}}
\expandafter\def\csname BookIndex@T0230\endcsname{zhenghehanzi@正合函子（exact functor）}
\expandafter\let\csname BookOriginal@trExactFunctor\endcsname\trExactFunctor
\renewcommand{\trExactFunctor}{\BookTerm{T0230}{\csname BookOriginal@trExactFunctor\endcsname}}
\expandafter\def\csname BookIndex@T0231\endcsname{shuangfanchou@双范畴（double category）}
\expandafter\let\csname BookOriginal@trDoubleCategory\endcsname\trDoubleCategory
\renewcommand{\trDoubleCategory}{\BookTerm{T0231}{\csname BookOriginal@trDoubleCategory\endcsname}}
\expandafter\def\csname BookIndex@T0232\endcsname{xunishuangfanchou@虚拟双范畴（virtual double category）}
\expandafter\let\csname BookOriginal@trVirtualDoubleCategory\endcsname\trVirtualDoubleCategory
\renewcommand{\trVirtualDoubleCategory}{\BookTerm{T0232}{\csname BookOriginal@trVirtualDoubleCategory\endcsname}}
\expandafter\def\csname BookIndex@T0233\endcsname{rouxingjiaquanjixian@柔性加权极限（flexible weighted limit）}
\expandafter\let\csname BookOriginal@trFlexibleWeightedLimit\endcsname\trFlexibleWeightedLimit
\renewcommand{\trFlexibleWeightedLimit}{\BookTerm{T0233}{\csname BookOriginal@trFlexibleWeightedLimit\endcsname}}
\expandafter\def\csname BookIndex@T0234\endcsname{yuzhoulunde@宇宙论的（cosmological）}
\expandafter\let\csname BookOriginal@trCosmological\endcsname\trCosmological
\renewcommand{\trCosmological}{\BookTerm{T0234}{\csname BookOriginal@trCosmological\endcsname}}
\expandafter\def\csname BookIndex@T0235\endcsname{yudikaerxianweihua@余笛卡尔纤维化（cocartesian fibration）}
\expandafter\let\csname BookOriginal@trCocartesianFibration\endcsname\trCocartesianFibration
\renewcommand{\trCocartesianFibration}{\BookTerm{T0235}{\csname BookOriginal@trCocartesianFibration\endcsname}}
\expandafter\def\csname BookIndex@T0236\endcsname{yuzhoulunqianrude@宇宙论嵌入的（cosmologically embedded）}
\expandafter\let\csname BookOriginal@trCosmologicallyEmbedded\endcsname\trCosmologicallyEmbedded
\renewcommand{\trCosmologicallyEmbedded}{\BookTerm{T0236}{\csname BookOriginal@trCosmologicallyEmbedded\endcsname}}
\expandafter\def\csname BookIndex@T0237\endcsname{shuangcedengjiaxianweihua@双侧等价纤维化（two-sided isofibration）}
\expandafter\let\csname BookOriginal@trTwoSidedIsofibration\endcsname\trTwoSidedIsofibration
\renewcommand{\trTwoSidedIsofibration}{\BookTerm{T0237}{\csname BookOriginal@trTwoSidedIsofibration\endcsname}}
\expandafter\def\csname BookIndex@T0238\endcsname{dikaerhanzi@笛卡尔函子（cartesian functor）}
\expandafter\let\csname BookOriginal@trCartesianFunctor\endcsname\trCartesianFunctor
\renewcommand{\trCartesianFunctor}{\BookTerm{T0238}{\csname BookOriginal@trCartesianFunctor\endcsname}}
\expandafter\def\csname BookIndex@T0239\endcsname{douhaokuadu@逗号跨度（comma span）}
\expandafter\let\csname BookOriginal@trCommaSpan\endcsname\trCommaSpan
\renewcommand{\trCommaSpan}{\BookTerm{T0239}{\csname BookOriginal@trCommaSpan\endcsname}}
\expandafter\def\csname BookIndex@T0240\endcsname{dianfande@典范的（canonical）}
\expandafter\let\csname BookOriginal@trCanonical\endcsname\trCanonical
\renewcommand{\trCanonical}{\BookTerm{T0240}{\csname BookOriginal@trCanonical\endcsname}}
\expandafter\def\csname BookIndex@T0241\endcsname{danzi@单子（monad）}
\expandafter\let\csname BookOriginal@trMonad\endcsname\trMonad
\renewcommand{\trMonad}{\BookTerm{T0241}{\csname BookOriginal@trMonad\endcsname}}
\expandafter\def\csname BookIndex@T0242\endcsname{yaobanfanchou@幺半范畴（monoidal category）}
\expandafter\let\csname BookOriginal@trMonoidalCategory\endcsname\trMonoidalCategory
\renewcommand{\trMonoidalCategory}{\BookTerm{T0242}{\csname BookOriginal@trMonoidalCategory\endcsname}}
\expandafter\def\csname BookIndex@T0243\endcsname{yaohuan@幺环（unital ring）}
\expandafter\let\csname BookOriginal@trUnitalRing\endcsname\trUnitalRing
\renewcommand{\trUnitalRing}{\BookTerm{T0243}{\csname BookOriginal@trUnitalRing\endcsname}}
\expandafter\def\csname BookIndex@T0244\endcsname{qianjiantouzhuangbei@前箭头装备（proarrow equipment）}
\expandafter\let\csname BookOriginal@trProarrowEquipment\endcsname\trProarrowEquipment
\renewcommand{\trProarrowEquipment}{\BookTerm{T0244}{\csname BookOriginal@trProarrowEquipment\endcsname}}
\expandafter\def\csname BookIndex@T0245\endcsname{weishuangfanchou@伪双范畴（pseudo double category）}
\expandafter\let\csname BookOriginal@trPseudoDoubleCategory\endcsname\trPseudoDoubleCategory
\renewcommand{\trPseudoDoubleCategory}{\BookTerm{T0245}{\csname BookOriginal@trPseudoDoubleCategory\endcsname}}
\expandafter\def\csname BookIndex@T0246\endcsname{tonglunyufanzhuanqi@同伦余反转器（homotopy coinverter）}
\expandafter\let\csname BookOriginal@trHomotopyCoinverter\endcsname\trHomotopyCoinverter
\renewcommand{\trHomotopyCoinverter}{\BookTerm{T0246}{\csname BookOriginal@trHomotopyCoinverter\endcsname}}
\expandafter\def\csname BookIndex@T0247\endcsname{kedade@可达的（accessible）}
\expandafter\let\csname BookOriginal@trAccessible\endcsname\trAccessible
\renewcommand{\trAccessible}{\BookTerm{T0247}{\csname BookOriginal@trAccessible\endcsname}}
\expandafter\def\csname BookIndex@T0248\endcsname{zhangliangji@张量积（tensor product）}
\expandafter\let\csname BookOriginal@trTensorProduct\endcsname\trTensorProduct
\renewcommand{\trTensorProduct}{\BookTerm{T0248}{\csname BookOriginal@trTensorProduct\endcsname}}
\expandafter\def\csname BookIndex@T0249\endcsname{biaoliangdexianzhi@标量的限制（restriction of scalars）}
\expandafter\let\csname BookOriginal@trRestrictionOfScalars\endcsname\trRestrictionOfScalars
\renewcommand{\trRestrictionOfScalars}{\BookTerm{T0249}{\csname BookOriginal@trRestrictionOfScalars\endcsname}}
\expandafter\def\csname BookIndex@T0250\endcsname{hanzi@∞-函子（\ensuremath{\infty}-functor）}
\expandafter\let\csname BookOriginal@trInfinityFunctor\endcsname\trInfinityFunctor
\renewcommand{\trInfinityFunctor}{\BookTerm{T0250}{\csname BookOriginal@trInfinityFunctor\endcsname}}
\expandafter\def\csname BookIndex@T0251\endcsname{1taishe@1-态射（1-morphism）}
\expandafter\let\csname BookOriginal@trOneMorphism\endcsname\trOneMorphism
\renewcommand{\trOneMorphism}{\BookTerm{T0251}{\csname BookOriginal@trOneMorphism\endcsname}}
\expandafter\def\csname BookIndex@T0252\endcsname{zhudiande@逐点的（pointwise）}
\expandafter\let\csname BookOriginal@trPointwise\endcsname\trPointwise
\renewcommand{\trPointwise}{\BookTerm{T0252}{\csname BookOriginal@trPointwise\endcsname}}
\expandafter\def\csname BookIndex@T0253\endcsname{zhudiankuozhang@逐点扩张（pointwise extension）}
\expandafter\let\csname BookOriginal@trPointwiseExtension\endcsname\trPointwiseExtension
\renewcommand{\trPointwiseExtension}{\BookTerm{T0253}{\csname BookOriginal@trPointwiseExtension\endcsname}}
\expandafter\def\csname BookIndex@T0254\endcsname{jiaquanyujixian@加权余极限（weighted colimit）}
\expandafter\let\csname BookOriginal@trWeightedColimit\endcsname\trWeightedColimit
\renewcommand{\trWeightedColimit}{\BookTerm{T0254}{\csname BookOriginal@trWeightedColimit\endcsname}}
\expandafter\def\csname BookIndex@T0255\endcsname{fanhanzi@泛函子（profunctor）}
\expandafter\let\csname BookOriginal@trProfunctor\endcsname\trProfunctor
\renewcommand{\trProfunctor}{\BookTerm{T0255}{\csname BookOriginal@trProfunctor\endcsname}}
\expandafter\def\csname BookIndex@T0256\endcsname{ruo2fanchou@弱 2-范畴（bicategory）}
\expandafter\let\csname BookOriginal@trBicategory\endcsname\trBicategory
\renewcommand{\trBicategory}{\BookTerm{T0256}{\csname BookOriginal@trBicategory\endcsname}}
\expandafter\def\csname BookIndex@T0257\endcsname{ruo2fanchoude@弱 2-范畴的（bicategorical）}
\expandafter\let\csname BookOriginal@trBicategorical\endcsname\trBicategorical
\renewcommand{\trBicategorical}{\BookTerm{T0257}{\csname BookOriginal@trBicategorical\endcsname}}
\expandafter\def\csname BookIndex@T0258\endcsname{suocunyingshe@锁存映射（latching map）}
\expandafter\let\csname BookOriginal@trLatchingMap\endcsname\trLatchingMap
\renewcommand{\trLatchingMap}{\BookTerm{T0258}{\csname BookOriginal@trLatchingMap\endcsname}}
\expandafter\def\csname BookIndex@T0259\endcsname{wutong@无痛（anodyne）}
\expandafter\let\csname BookOriginal@trAnodyne\endcsname\trAnodyne
\renewcommand{\trAnodyne}{\BookTerm{T0259}{\csname BookOriginal@trAnodyne\endcsname}}
\expandafter\def\csname BookIndex@T0260\endcsname{wutongkuozhang@无痛扩张（anodyne extension）}
\expandafter\let\csname BookOriginal@trAnodyneExtension\endcsname\trAnodyneExtension
\renewcommand{\trAnodyneExtension}{\BookTerm{T0260}{\csname BookOriginal@trAnodyneExtension\endcsname}}
\expandafter\def\csname BookIndex@T0261\endcsname{neiwutong@内无痛（inner anodyne）}
\expandafter\let\csname BookOriginal@trInnerAnodyne\endcsname\trInnerAnodyne
\renewcommand{\trInnerAnodyne}{\BookTerm{T0261}{\csname BookOriginal@trInnerAnodyne\endcsname}}
\expandafter\def\csname BookIndex@T0262\endcsname{zhengbaohan@整包含（entire inclusion）}
\expandafter\let\csname BookOriginal@trEntireInclusion\endcsname\trEntireInclusion
\renewcommand{\trEntireInclusion}{\BookTerm{T0262}{\csname BookOriginal@trEntireInclusion\endcsname}}
\expandafter\def\csname BookIndex@T0263\endcsname{lianjie@联接（join）}
\expandafter\let\csname BookOriginal@trJoin\endcsname\trJoin
\renewcommand{\trJoin}{\BookTerm{T0263}{\csname BookOriginal@trJoin\endcsname}}
\expandafter\def\csname BookIndex@T0264\endcsname{panglianjie@胖联接（fat join）}
\expandafter\let\csname BookOriginal@trFatJoin\endcsname\trFatJoin
\renewcommand{\trFatJoin}{\BookTerm{T0264}{\csname BookOriginal@trFatJoin\endcsname}}
\expandafter\def\csname BookIndex@T0265\endcsname{jueduiyoutisheng@绝对右提升（absolute right lifting）}
\expandafter\let\csname BookOriginal@trAbsoluteRightLifting\endcsname\trAbsoluteRightLifting
\renewcommand{\trAbsoluteRightLifting}{\BookTerm{T0265}{\csname BookOriginal@trAbsoluteRightLifting\endcsname}}
\expandafter\def\csname BookIndex@T0266\endcsname{jueduizuotisheng@绝对左提升（absolute left lifting）}
\expandafter\let\csname BookOriginal@trAbsoluteLeftLifting\endcsname\trAbsoluteLeftLifting
\renewcommand{\trAbsoluteLeftLifting}{\BookTerm{T0266}{\csname BookOriginal@trAbsoluteLeftLifting\endcsname}}
\expandafter\def\csname BookIndex@T0267\endcsname{neixianweihua@内纤维化（inner fibration）}
\expandafter\let\csname BookOriginal@trInnerFibration\endcsname\trInnerFibration
\renewcommand{\trInnerFibration}{\BookTerm{T0267}{\csname BookOriginal@trInnerFibration\endcsname}}
\expandafter\def\csname BookIndex@T0268\endcsname{dikaerjiantou@笛卡尔箭头（cartesian arrow）}
\expandafter\let\csname BookOriginal@trCartesianArrow\endcsname\trCartesianArrow
\renewcommand{\trCartesianArrow}{\BookTerm{T0268}{\csname BookOriginal@trCartesianArrow\endcsname}}
\expandafter\def\csname BookIndex@T0269\endcsname{yudikaerjiantou@余笛卡尔箭头（cocartesian arrow）}
\expandafter\let\csname BookOriginal@trCocartesianArrow\endcsname\trCocartesianArrow
\renewcommand{\trCocartesianArrow}{\BookTerm{T0269}{\csname BookOriginal@trCocartesianArrow\endcsname}}
\expandafter\def\csname BookIndex@T0270\endcsname{youxianweihua@右纤维化（right fibration）}
\expandafter\let\csname BookOriginal@trRightFibration\endcsname\trRightFibration
\renewcommand{\trRightFibration}{\BookTerm{T0270}{\csname BookOriginal@trRightFibration\endcsname}}
\expandafter\def\csname BookIndex@T0271\endcsname{zuoxianweihua@左纤维化（left fibration）}
\expandafter\let\csname BookOriginal@trLeftFibration\endcsname\trLeftFibration
\renewcommand{\trLeftFibration}{\BookTerm{T0271}{\csname BookOriginal@trLeftFibration\endcsname}}
\expandafter\def\csname BookIndex@T0272\endcsname{duiying@对应（correspondence）}
\expandafter\let\csname BookOriginal@trCorrespondence\endcsname\trCorrespondence
\renewcommand{\trCorrespondence}{\BookTerm{T0272}{\csname BookOriginal@trCorrespondence\endcsname}}
\expandafter\def\csname BookIndex@T0273\endcsname{pintie@拼贴（collage）}
\expandafter\let\csname BookOriginal@trCollage\endcsname\trCollage
\renewcommand{\trCollage}{\BookTerm{T0273}{\csname BookOriginal@trCollage\endcsname}}
\expandafter\def\csname BookIndex@T0274\endcsname{fansongchiyujixian@反松弛余极限（oplax colimit）}
\expandafter\let\csname BookOriginal@trOplaxColimit\endcsname\trOplaxColimit
\renewcommand{\trOplaxColimit}{\BookTerm{T0274}{\csname BookOriginal@trOplaxColimit\endcsname}}
\expandafter\def\csname BookIndex@T0275\endcsname{shuangxianweihua@双纤维化（bifibration）}
\expandafter\let\csname BookOriginal@trBifibration\endcsname\trBifibration
\renewcommand{\trBifibration}{\BookTerm{T0275}{\csname BookOriginal@trBifibration\endcsname}}
\expandafter\def\csname BookIndex@T0276\endcsname{fanchouquanji@范畴圈积（categorial wreath product）}
\expandafter\let\csname BookOriginal@trWreathProduct\endcsname\trWreathProduct
\renewcommand{\trWreathProduct}{\BookTerm{T0276}{\csname BookOriginal@trWreathProduct\endcsname}}
\expandafter\def\csname BookIndex@T0277\endcsname{rezkduixiang@Rezk 对象（Rezk object）}
\expandafter\let\csname BookOriginal@trRezkObject\endcsname\trRezkObject
\renewcommand{\trRezkObject}{\BookTerm{T0277}{\csname BookOriginal@trRezkObject\endcsname}}
\expandafter\def\csname BookIndex@T0278\endcsname{rezkruodengjia@Rezk 弱等价（Rezk weak equivalence）}
\expandafter\let\csname BookOriginal@trRezkWeakEquivalence\endcsname\trRezkWeakEquivalence
\renewcommand{\trRezkWeakEquivalence}{\BookTerm{T0278}{\csname BookOriginal@trRezkWeakEquivalence\endcsname}}
\expandafter\def\csname BookIndex@T0279\endcsname{rezkdengjiaxianweihua@Rezk 等价纤维化（Rezk isofibration）}
\expandafter\let\csname BookOriginal@trRezkIsofibration\endcsname\trRezkIsofibration
\renewcommand{\trRezkIsofibration}{\BookTerm{T0279}{\csname BookOriginal@trRezkIsofibration\endcsname}}
\expandafter\def\csname BookIndex@T0280\endcsname{segalduixiang@Segal 对象（Segal object）}
\expandafter\let\csname BookOriginal@trSegalObject\endcsname\trSegalObject
\renewcommand{\trSegalObject}{\BookTerm{T0280}{\csname BookOriginal@trSegalObject\endcsname}}
\expandafter\def\csname BookIndex@T0281\endcsname{nzhongwanquansegalkongjian@n-重完全 Segal 空间（n-fold complete Segal space）}
\expandafter\let\csname BookOriginal@trNFoldCompleteSegalSpace\endcsname\trNFoldCompleteSegalSpace
\renewcommand{\trNFoldCompleteSegalSpace}{\BookTerm{T0281}{\csname BookOriginal@trNFoldCompleteSegalSpace\endcsname}}
\expandafter\def\csname BookIndex@T0282\endcsname{cisinskimoxingfanchou@Cisinski 模型范畴（Cisinski model category）}
\expandafter\let\csname BookOriginal@trCisinskiModelCategory\endcsname\trCisinskiModelCategory
\renewcommand{\trCisinskiModelCategory}{\BookTerm{T0282}{\csname BookOriginal@trCisinskiModelCategory\endcsname}}
\expandafter\def\csname BookIndex@T0283\endcsname{cisinskimoxingjiegou@Cisinski 模型结构（Cisinski model structure）}
\expandafter\let\csname BookOriginal@trCisinskiModelStructure\endcsname\trCisinskiModelStructure
\renewcommand{\trCisinskiModelStructure}{\BookTerm{T0283}{\csname BookOriginal@trCisinskiModelStructure\endcsname}}
\expandafter\def\csname BookIndex@T0284\endcsname{heou@核偶（kernel pair）}
\expandafter\let\csname BookOriginal@trKernelPair\endcsname\trKernelPair
\renewcommand{\trKernelPair}{\BookTerm{T0284}{\csname BookOriginal@trKernelPair\endcsname}}
\expandafter\def\csname BookIndex@T0285\endcsname{jubuhuazi@局部化子（localizer）}
\expandafter\let\csname BookOriginal@trLocalizer\endcsname\trLocalizer
\renewcommand{\trLocalizer}{\BookTerm{T0285}{\csname BookOriginal@trLocalizer\endcsname}}
\expandafter\def\csname BookIndex@T0286\endcsname{zuobousfieldjubuhua@左 Bousfield 局部化（left Bousfield localization）}
\expandafter\let\csname BookOriginal@trBousfieldLocalization\endcsname\trBousfieldLocalization
\renewcommand{\trBousfieldLocalization}{\BookTerm{T0286}{\csname BookOriginal@trBousfieldLocalization\endcsname}}
\expandafter\def\csname BookIndex@T0287\endcsname{neishemoxingjiegou@内射模型结构（injective model structure）}
\expandafter\let\csname BookOriginal@trInjectiveModelStructure\endcsname\trInjectiveModelStructure
\renewcommand{\trInjectiveModelStructure}{\BookTerm{T0287}{\csname BookOriginal@trInjectiveModelStructure\endcsname}}
\expandafter\def\csname BookIndex@T0288\endcsname{grothendiecktapusi@Grothendieck 拓扑斯（Grothendieck topos）}
\expandafter\let\csname BookOriginal@trGrothendieckTopos\endcsname\trGrothendieckTopos
\renewcommand{\trGrothendieckTopos}{\BookTerm{T0288}{\csname BookOriginal@trGrothendieckTopos\endcsname}}
\expandafter\def\csname BookIndex@T0289\endcsname{qiumian@球面（sphere）}
\expandafter\let\csname BookOriginal@trSphere\endcsname\trSphere
\renewcommand{\trSphere}{\BookTerm{T0289}{\csname BookOriginal@trSphere\endcsname}}
\expandafter\def\csname BookIndex@T0290\endcsname{minsumoxingjiegou@民俗模型结构（folk model structure）}
\expandafter\let\csname BookOriginal@trFolkModelStructure\endcsname\trFolkModelStructure
\renewcommand{\trFolkModelStructure}{\BookTerm{T0290}{\csname BookOriginal@trFolkModelStructure\endcsname}}
\expandafter\def\csname BookIndex@T0291\endcsname{tonggoudouhaofanchou@同构逗号范畴（iso-comma category）}
\expandafter\let\csname BookOriginal@trIsoCommaCategory\endcsname\trIsoCommaCategory
\renewcommand{\trIsoCommaCategory}{\BookTerm{T0291}{\csname BookOriginal@trIsoCommaCategory\endcsname}}
\expandafter\def\csname BookIndex@T0292\endcsname{leibnizji@Leibniz 积（Leibniz product）}
\expandafter\let\csname BookOriginal@trLeibnizProduct\endcsname\trLeibnizProduct
\renewcommand{\trLeibnizProduct}{\BookTerm{T0292}{\csname BookOriginal@trLeibnizProduct\endcsname}}
\expandafter\def\csname BookIndex@T0293\endcsname{yaobanchengji@幺半乘积（monoidal product）}
\expandafter\let\csname BookOriginal@trMonoidalProduct\endcsname\trMonoidalProduct
\renewcommand{\trMonoidalProduct}{\BookTerm{T0293}{\csname BookOriginal@trMonoidalProduct\endcsname}}
\expandafter\def\csname BookIndex@T0294\endcsname{tongtai@同态（homomorphism）}
\expandafter\let\csname BookOriginal@trHomomorphism\endcsname\trHomomorphism
\renewcommand{\trHomomorphism}{\BookTerm{T0294}{\csname BookOriginal@trHomomorphism\endcsname}}
\expandafter\def\csname BookIndex@T0295\endcsname{danchunjisuanzi@单纯计算子（simplicial computad）}
\expandafter\let\csname BookOriginal@trSimplicialComputad\endcsname\trSimplicialComputad
\renewcommand{\trSimplicialComputad}{\BookTerm{T0295}{\csname BookOriginal@trSimplicialComputad\endcsname}}
\expandafter\def\csname BookIndex@T0296\endcsname{danchunyingshe@单纯映射（simplicial map）}
\expandafter\let\csname BookOriginal@trSimplicialMap\endcsname\trSimplicialMap
\renewcommand{\trSimplicialMap}{\BookTerm{T0296}{\csname BookOriginal@trSimplicialMap\endcsname}}
\expandafter\def\csname BookIndex@T0297\endcsname{douhaomo@逗号模（comma module）}
\expandafter\let\csname BookOriginal@trCommaModule\endcsname\trCommaModule
\renewcommand{\trCommaModule}{\BookTerm{T0297}{\csname BookOriginal@trCommaModule\endcsname}}
\expandafter\def\csname BookIndex@T0298\endcsname{douhaofangkuai@逗号方块（comma square）}
\expandafter\let\csname BookOriginal@trCommaSquare\endcsname\trCommaSquare
\renewcommand{\trCommaSquare}{\BookTerm{T0298}{\csname BookOriginal@trCommaSquare\endcsname}}
\expandafter\def\csname BookIndex@T0299\endcsname{banlu@伴侣（companion）}
\expandafter\let\csname BookOriginal@trCompanion\endcsname\trCompanion
\renewcommand{\trCompanion}{\BookTerm{T0299}{\csname BookOriginal@trCompanion\endcsname}}
\expandafter\def\csname BookIndex@T0300\endcsname{gongou@共偶（conjoint）}
\expandafter\let\csname BookOriginal@trConjoint\endcsname\trConjoint
\renewcommand{\trConjoint}{\BookTerm{T0300}{\csname BookOriginal@trConjoint\endcsname}}
\expandafter\def\csname BookIndex@T0301\endcsname{qianfuhe@前复合（pre-composition）}
\expandafter\let\csname BookOriginal@trPreComposition\endcsname\trPreComposition
\renewcommand{\trPreComposition}{\BookTerm{T0301}{\csname BookOriginal@trPreComposition\endcsname}}
\expandafter\def\csname BookIndex@T0302\endcsname{xianweihuabansui@纤维化伴随（fibered adjunction）}
\expandafter\let\csname BookOriginal@trFiberedAdjunction\endcsname\trFiberedAdjunction
\renewcommand{\trFiberedAdjunction}{\BookTerm{T0302}{\csname BookOriginal@trFiberedAdjunction\endcsname}}
\expandafter\def\csname BookIndex@T0303\endcsname{lisandikaerxianweihua@离散笛卡尔纤维化（discrete cartesian fibration）}
\expandafter\let\csname BookOriginal@trDiscreteCartesianFibration\endcsname\trDiscreteCartesianFibration
\renewcommand{\trDiscreteCartesianFibration}{\BookTerm{T0303}{\csname BookOriginal@trDiscreteCartesianFibration\endcsname}}
\expandafter\def\csname BookIndex@T0304\endcsname{lisanyudikaerxianweihua@离散余笛卡尔纤维化（discrete cocartesian fibration）}
\expandafter\let\csname BookOriginal@trDiscreteCocartesianFibration\endcsname\trDiscreteCocartesianFibration
\renewcommand{\trDiscreteCocartesianFibration}{\BookTerm{T0304}{\csname BookOriginal@trDiscreteCocartesianFibration\endcsname}}
\expandafter\def\csname BookIndex@T0305\endcsname{fanchou@∞-范畴（\ensuremath{\infty} categories）}
\expandafter\let\csname BookOriginal@trInfinityCategories\endcsname\trInfinityCategories
\renewcommand{\trInfinityCategories}{\BookTerm{T0305}{\csname BookOriginal@trInfinityCategories\endcsname}}
\makeatother
